\documentclass[11pt,a4paper]{article}

% ========================================
% PACKAGES
% ========================================
\usepackage[utf8]{inputenc}
\usepackage[T1]{fontenc}
\usepackage{lmodern}
\usepackage[english]{babel}
\usepackage[margin=2.5cm]{geometry}

% Mathematics
\usepackage{amsmath, amssymb, amsthm}
\usepackage{mathtools}

% Graphics and colors
\usepackage{graphicx}
\usepackage{xcolor}
\usepackage{tikz}
\usetikzlibrary{positioning, arrows.meta, shapes.geometric, calc}
\usepackage{pgfplots}
\pgfplotsset{compat=1.18}
\usepackage[most]{tcolorbox}

% Lists and formatting
\usepackage{enumitem}
\usepackage{parskip}
\usepackage{fancyhdr}
\usepackage{multirow}

% Code listings
\usepackage{listings}
\usepackage{algorithm}
\usepackage{algorithmic}

% Hyperlinks
\usepackage{hyperref}
\hypersetup{
    colorlinks=true,
    linkcolor=blue,
    urlcolor=blue,
    citecolor=blue
}

% ========================================
% THEOREM ENVIRONMENTS
% ========================================
\theoremstyle{definition}
\newtheorem{definition}{Definition}[section]
\newtheorem{example}{Example}[section]
\newtheorem{exercise}{Exercise}[section]

\theoremstyle{plain}
\newtheorem{theorem}{Theorem}[section]
\newtheorem{lemma}[theorem]{Lemma}
\newtheorem{proposition}[theorem]{Proposition}
\newtheorem{corollary}[theorem]{Corollary}

\theoremstyle{remark}
\newtheorem*{remark}{Remark}
\newtheorem*{observation}{Observation}

% ========================================
% CUSTOM COMMANDS
% ========================================
\newcommand{\R}{\mathbb{R}}
\newcommand{\N}{\mathbb{N}}
\newcommand{\Z}{\mathbb{Z}}
\newcommand{\Q}{\mathbb{Q}}
\newcommand{\C}{\mathbb{C}}

% ========================================
% TABLE OF CONTENTS DEPTH
% ========================================
\setcounter{tocdepth}{2}

% ========================================
% HEADER AND FOOTER
% ========================================
\setlength{\headheight}{14pt}
\pagestyle{fancy}
\fancyhf{}
\lhead{\leftmark}
\rhead{SL}
\cfoot{\thepage}
\renewcommand{\sectionmark}[1]{\markboth{#1}{}}

% ========================================
% DOCUMENT INFORMATION
% ========================================
\title{\textbf{Statistical Learning}\\
\large Artificial Intelligence}
\author{Jacopo Parretti}
\date{II Semester 2025-2026}

% ========================================
% DOCUMENT
% ========================================
\begin{document}

\maketitle
\tableofcontents
\newpage


% ========================================
\part{Introduction to Statistical Learning}
% ========================================

\section{Motivations}

\subsection{What is Statistical Learning?}

\begin{tcolorbox}[colback=blue!5!white,colframe=blue!75!black,title=\textbf{Statistical Learning}]
Statistical learning refers to a vast set of \textbf{tools for understanding data}. These tools enable us to model, predict, and draw inferences from complex datasets across a wide range of application domains.
\end{tcolorbox}

The central goal is to \textbf{understand correlations between data}. Depending on what we want to predict and what kind of output we have, this takes three main forms:

\begin{itemize}
    \item \textbf{Regression}: predicting a \textit{continuous} (quantitative) output variable.\\
    \textit{Example}: predicting wage as a function of age, year of experience, and education level.
    \item \textbf{Classification}: predicting a \textit{discrete} (categorical) output variable.\\
    \textit{Example}: predicting the direction of S\&P 500 daily movement (up/down) from previous days' percentage changes.
    \item \textbf{Clustering}: \textit{grouping} data points based on observed characteristics, with no output label available.\\
    \textit{Example}: grouping cell lines by gene expression profiles (unsupervised — no cancer type label used during grouping).
\end{itemize}

\subsection{A Brief History}

\begin{itemize}
    \item \textbf{Early 1800s} — Legendre and Gauss: least squares method, linear regression.
    \item \textbf{Early 1936} — Fisher: Linear Discriminant Analysis, classification.
    \item \textbf{Mid 1980s} — Breiman, Friedman, Olshen, Stone: classification and regression trees; tree-based methods and cross-validation for model selection.
    \item \textbf{Present} — Machine Learning, Data Science, Deep Learning.
\end{itemize}

\subsection{Why Study Statistical Learning?}

Statistical learning methods are key assets for a wide range of disciplines:
\begin{itemize}
    \item Artificial Intelligence, Machine Learning, Economics, Business, Social Sciences, Medicine, Biotechnology, Public Policies
    \item Finance, Manufacturing, Pharmaceutical, Logistics, Entertainment, Politics, Tourism
\end{itemize}

Three core principles motivate this course:
\begin{enumerate}
    \item Statistical learning methods should \textbf{not} be treated as black boxes: to select the best method you must understand how they work.
    \item You do not need to re-implement state-of-the-art methods, but you must \textbf{understand them and know when and how to use them}.
    \item The goal is to develop competence in the \textbf{application of statistical learning methods to real-world problems} (data science).
\end{enumerate}

\subsection{Real-World Applications}

Two representative application areas illustrate the scope of the discipline.

\subsubsection{District Heating Networks}

\textbf{Goal}: predict heating load in the next 48 hours.
\begin{itemize}
    \item \textbf{Data}: multivariate time series (heat demand of previous days, external temperature, humidity, \ldots)
    \item \textbf{Model}: a set of linear autoregressive models
    \item \textbf{Output}: load demand forecast for the next 48 hours
\end{itemize}
This is a regression problem. Evaluation metrics used in practice include RMSE and MAPE across different regimes (weekdays, weekends, holidays).

\subsubsection{Environmental Monitoring — Autonomous Water Drones}

\textbf{Goal}: recognize what a drone is doing from its sensor stream (activity recognition).
\begin{itemize}
    \item \textbf{Data}: multivariate time series from GPS, engine control, battery, electrical conductivity, dissolved oxygen, temperature sensors
    \item \textbf{Task}: find time intervals representing important situations (e.g., drone moving against current)
    \item \textbf{Approach}: unsupervised learning — K-means, Affinity Propagation, Gaussian Mixture Models, Hidden Markov Models
\end{itemize}
Key result: it is possible to \textbf{segment time series} and identify meaningful clusters without knowing the labels. For upstream/downstream navigation, K-means identifies average direction of movement, average voltage to motors, and standard deviation of battery voltage as the most relevant features.


\section{Key Elements of Statistical Learning}

\subsection{The Core Idea}

\begin{tcolorbox}[colback=blue!5!white,colframe=blue!75!black,title=\textbf{Core idea}]
Statistical learning is about \textbf{automatically building models starting from data}. In all cases, what we learn is a \textbf{function}.
\end{tcolorbox}

There are three main types of learning, distinguished by what information is available during training:

\begin{itemize}
    \item \textbf{Supervised learning}: $y = f(x)$ — given pairs $(x, y)$, learn $f(\cdot)$ [function approximation].
    \item \textbf{Unsupervised learning}: $f(x)$ — given only $x$, learn $f(\cdot)$ as a compact representation of $x$ [clustering].
    \item \textbf{Reinforcement learning}: $y = f(x),\, z$ — given $(x, z)$, learn $f(\cdot)$ that generates $y$, where $x$ is the state and $z$ is the reward signal [sequential decision making].
\end{itemize}

\subsection{Supervised Learning}

In supervised learning we have:
\begin{itemize}
    \item A \textbf{training set} of labeled examples: known inputs $x$ and corresponding labels $t$.
    \item A \textbf{model} $y(x)$ that maps inputs to predicted outputs.
    \item A \textbf{test set} of new, unseen inputs: the model is applied to evaluate generalization.
\end{itemize}

\textit{Example}: recognizing hand-written digits. Training data consists of images $x$ and digit labels $t$. The model learns a mapping $y(x)$; at test time it predicts the digit for new images.

\subsection{Unsupervised Learning}

In unsupervised learning there are \textbf{no labels}. The goal is to build a model that \textbf{groups similar} data points — by columns, rows, or both (bi-clustering).

\textit{Example}: gene expression data. Rows are genes, columns are cell lines, each entry is the expression level. No cancer-type label is used; the task is to discover structure in the data.

\subsection{Reinforcement Learning}

Reinforcement learning is a \textbf{control problem with sequential decision making}:
\begin{itemize}
    \item A \textbf{reward signal} drives the agent's actions (e.g., high reward when reaching a target).
    \item The agent must \textbf{interact with the environment} to learn.
    \item It uses data-driven statistical techniques but the focus is on \textbf{control}, not data analysis.
\end{itemize}

\subsection{Basic Concepts and Terminology}

\begin{itemize}
    \item \textbf{Input variables} (independent variables, predictors, features): denoted $X$ when scalar, $X_1, X_2, \ldots, X_p$ when there are $p$ predictors.
    \item \textbf{Output variable} (response, dependent variable): denoted $Y$.
    \item \textbf{Quantitative variables}: numerical values (temperature, height, income, \ldots).
    \item \textbf{Qualitative variables} (categorical): values from one of $K$ classes (gender, brand, \ldots).
    \item \textbf{Regression}: quantitative output.
    \item \textbf{Classification}: qualitative output.
\end{itemize}


\section{Why Estimate \texorpdfstring{$f$}{f}?}

\subsection{The General Setup}

Assume $Y = f(X) + \epsilon$, where $f$ is an unknown function of $X = (X_1, X_2, \ldots, X_p)$ and $\epsilon$ is a \textbf{random error term} independent of $X$. Statistical learning is concerned with estimating $f$.

There are two main motivations for estimating $f$: \textbf{prediction} and \textbf{inference}.

\subsection{Prediction}

When inputs $X$ are easy to observe but output $Y$ is difficult to measure directly, we build a model $\hat{f}(\cdot)$ to compute:
\[
    \hat{Y} = \hat{f}(X)
\]
In prediction, the exact form of $\hat{f}$ is not of primary interest — it acts as a \textbf{black box} as long as predictions are accurate.

\subsubsection{Reducible and Irreducible Error}

The accuracy of $\hat{Y}$ as a prediction of $Y$ depends on two quantities:

\begin{itemize}
    \item \textbf{Reducible error}: error that arises because $\hat{f}$ is not a perfect estimate of $f$. This can in principle be reduced by using better models or more data.
    \item \textbf{Irreducible error}: variability of $\epsilon$ (the random error), which is independent of $X$ and cannot be eliminated regardless of how well we estimate $f$.
\end{itemize}

Formally, the expected squared prediction error decomposes as:
\[
    E\!\left[(Y - \hat{Y})^2\right]
    = E\!\left[(f(X) - \hat{f}(X))^2\right] + \mathrm{Var}(\epsilon)
    = \underbrace{E\!\left[(f(X) - \hat{f}(X))^2\right]}_{\text{Reducible error}} + \underbrace{\mathrm{Var}(\epsilon)}_{\text{Irreducible error}}
\]

where $E(X)$ is the expected value (population average) of $X$ and $\mathrm{Var}(X)$ is its variance.

\subsection{Inference}

When the goal is to \textbf{understand how $Y$ changes as $X$ changes}, prediction accuracy is secondary — we need to know the \textbf{form} of $f(\cdot)$, so $\hat{f}$ cannot be a black box.

Typical questions in inference:
\begin{itemize}
    \item Which independent variables are \textbf{associated} with the response? (key for interpretability)
    \item What is the \textbf{relationship} between the response and each predictor: positive or negative effect?
    \item Can the relationship between $Y$ and $X$ be captured by a \textbf{linear model}? Linear models are simpler and easier to use but have limited representational power.
\end{itemize}

\subsection{Estimating \texorpdfstring{$f$}{f}: General Setup}

Given:
\begin{itemize}
    \item $n$ training data-points, $p$ predictors
    \item $x_{ij}$: value of observation $i$ for predictor $j$, where $i = 1, \ldots, n$ and $j = 1, \ldots, p$
    \item $y_i$: response value for observation $i$
    \item Training set: $(x_1, y_1), (x_2, y_2), \ldots, (x_n, y_n)$ where $x_i = (x_{i1}, x_{i2}, \ldots, x_{ip})^T$
\end{itemize}

\textbf{Goal}: find $\hat{f}(\cdot)$ such that $Y \approx \hat{f}(X)$ for any observation $(X, Y)$.


\section{How Do We Estimate \texorpdfstring{$f$}{f}?}

\subsection{Parametric Methods}

Parametric methods follow a \textbf{two-step procedure}:

\begin{enumerate}
    \item \textbf{Assume a form} for $f(\cdot)$. For example, assume $f$ is \textit{linear}:
    \[
        f(X) = \beta_0 + \beta_1 X_1 + \beta_2 X_2 + \cdots + \beta_p X_p
    \]
    A model is then completely identified by $p+1$ parameters $\beta_0, \beta_1, \ldots, \beta_p$.
    \item \textbf{Fit or train} the model using the training data. For a linear model this means estimating $\beta_0, \beta_1, \ldots, \beta_p$ such that
    \[
        Y \approx \beta_0 + \beta_1 X_1 + \beta_2 X_2 + \cdots + \beta_p X_p
    \]
    The most common approach is \textbf{Least Squares}.
\end{enumerate}

The key point: by assuming a parametric form we reduce the problem of estimating an arbitrary function $f$ to estimating a finite set of parameters.

\subsection{Non-Parametric Methods}

Non-parametric methods make \textbf{no assumption on the form of $f(\cdot)$} and instead try to fit $f$ as closely as possible to the training data.

\begin{itemize}
    \item \textbf{Pros}: can approximate the training data very well, no structural assumption.
    \item \textbf{Cons}: require a very large number of observations; can result in \textbf{overfitting}.
\end{itemize}

A typical example is thin-plate splines. A smooth spline fits the data well while generalizing; a rough spline interpolates the training points almost exactly but overfits.

\begin{tcolorbox}[colback=orange!5!white,colframe=orange!75!black,title=\textbf{Overfitting}]
Overfitting occurs when a model fits the training data too closely, capturing noise rather than the true underlying pattern. The model performs well on training data but poorly on new, unseen data.
\end{tcolorbox}


\section{Prediction Accuracy vs.\ Model Interpretability}

There is a fundamental \textbf{trade-off between flexibility and interpretability} in statistical learning methods. Along the flexibility axis, from most interpretable to most flexible:

\begin{enumerate}
    \item \textbf{Subset selection, Lasso} — high interpretability, low flexibility
    \item \textbf{Least Squares} — moderate interpretability
    \item \textbf{Generalized Additive Models, Trees} — moderate
    \item \textbf{Bagging, Boosting} — low interpretability, high flexibility
    \item \textbf{Support Vector Machines} — low interpretability, high flexibility
\end{enumerate}

More flexible methods tend to achieve better \textbf{prediction accuracy} on complex data but at the cost of \textbf{interpretability}. When inference is the goal, simpler models are preferred; when prediction accuracy is paramount and the model is used as a black box, more flexible methods may be appropriate.


\section{Assessing Model Accuracy}

\subsection{The Problem of Model Selection}

No method outperforms all others on all datasets. Selecting the best approach for a given problem is a challenging and crucial task. The key question is: how do we measure \textbf{quality of fit} — how far is the predicted response from the true response?

\subsection{Mean Squared Error}

For regression problems, the standard measure is the \textbf{Mean Squared Error (MSE)}:
\[
    \mathrm{MSE} = \frac{1}{n} \sum_{i=1}^{n} \bigl(y_i - \hat{f}(x_i)\bigr)^2
\]

\begin{itemize}
    \item \textbf{Training MSE}: MSE computed on the training data. Minimizing training MSE is how models are typically fit.
    \item \textbf{Test MSE}: MSE computed on new, unseen data. This is what we actually care about.
\end{itemize}

\begin{tcolorbox}[colback=orange!5!white,colframe=orange!75!black]
There is \textbf{no guarantee} that the method with minimum training MSE will achieve minimum test MSE. A model that fits the training data too well (low training MSE) often performs worse on test data (high test MSE) — overfitting.
\end{tcolorbox}

\subsection{The Bias-Variance Trade-off}

The expected test MSE at a given point $x_0$ can always be decomposed into three components:

\[
    E\!\left[\bigl(y_0 - \hat{f}(x_0)\bigr)^2\right]
    = \mathrm{Var}\!\left(\hat{f}(x_0)\right) + \left[\mathrm{Bias}\!\left(\hat{f}(x_0)\right)\right]^2 + \mathrm{Var}(\epsilon)
\]

\begin{itemize}
    \item \textbf{Variance} $\mathrm{Var}(\hat{f}(x_0))$: how much $\hat{f}$ would change if estimated on a different training set. More flexible models have \textbf{higher variance}.
    \item \textbf{Bias} $[\mathrm{Bias}(\hat{f}(x_0))]^2$: error introduced by the assumptions made on the model. Simpler models (e.g., linear) have \textbf{higher bias}.
    \item \textbf{Irreducible error} $\mathrm{Var}(\epsilon)$: irreducible noise, independent of the model.
\end{itemize}

\begin{tcolorbox}[colback=blue!5!white,colframe=blue!75!black,title=\textbf{Bias-Variance Trade-off}]
As model flexibility increases, \textbf{variance increases} and \textbf{bias decreases}. The optimal model minimizes the sum of variance and squared bias. This trade-off is unavoidable — there is no universally best model.
\end{tcolorbox}

Different datasets exhibit different bias-variance trade-offs (different rates of change). In practice, we cannot compute the test MSE directly, but we can \textbf{estimate} it — for example, using \textbf{cross-validation}.


% ========================================
\part{Linear Regression}
% ========================================

Linear regression is a useful tool for predicting a \textbf{quantitative} response. It is a simple yet powerful approach that introduces the supervised learning perspective, and serves as a jumping-off point for more modern statistical learning methods — many of which can be seen as generalizations or extensions of linear regression.

\section{Simple Linear Regression}

\subsection{The Model}

Simple linear regression predicts a quantitative response $Y$ on the basis of a \textbf{single predictor} $X$:
\[
    Y \approx \beta_0 + \beta_1 X
\]
We read $\approx$ as ``is approximately modelled as''. The \textbf{intercept} $\beta_0$ and the \textbf{slope} $\beta_1$ are the model \textbf{coefficients} (or parameters). Once we have estimated $\hat{\beta}_0$ and $\hat{\beta}_1$ from training data, predictions are computed as:
\[
    \hat{y} = \hat{\beta}_0 + \hat{\beta}_1 x
\]

\textit{Running example}: the Advertising dataset contains sales of a product in 200 different markets, together with advertising budgets for TV, radio, and newspaper. We can regress sales $Y$ on TV spending $X$:
\[
    \text{sales} \approx \beta_0 + \beta_1 \times \text{TV}
\]

\subsection{Estimating the Coefficients via Least Squares}

$\beta_0$ and $\beta_1$ are unknown. Starting from $n$ observation pairs $(x_1, y_1), \ldots, (x_n, y_n)$, we define the \textbf{residual} for observation $i$ as:
\[
    e_i = y_i - \hat{y}_i
\]
the difference between the $i$-th observed and predicted response. The \textbf{Residual Sum of Squares (RSS)} is:
\[
    \mathrm{RSS} = e_1^2 + \cdots + e_n^2
    = \sum_{i=1}^{n}(y_i - \hat{\beta}_0 - \hat{\beta}_1 x_i)^2
\]
The least squares estimates minimize the RSS. The closed-form solution is:
\[
    \hat{\beta}_1 = \frac{\sum_{i=1}^{n}(x_i - \bar{x})(y_i - \bar{y})}{\sum_{i=1}^{n}(x_i - \bar{x})^2}, \qquad
    \hat{\beta}_0 = \bar{y} - \hat{\beta}_1 \bar{x}
\]
where $\bar{x} = \frac{1}{n}\sum_{i=1}^n x_i$ and $\bar{y} = \frac{1}{n}\sum_{i=1}^n y_i$ are the sample means.

For the TV advertising example: $\hat{\beta}_0 = 7.03$, $\hat{\beta}_1 = 0.0475$.

\subsection{The Population Regression Line}

\begin{tcolorbox}[colback=blue!5!white,colframe=blue!75!black,title=\textbf{Population Regression Line}]
When $f$ can be approximated by a linear function, the \textbf{population regression line} is the model
\[
    Y = \beta_0 + \beta_1 X + \epsilon
\]
that defines the best linear approximation to the \textbf{true relationship} between $X$ and $Y$.
\begin{itemize}
    \item $\beta_0$: the intercept — expected value of $Y$ when $X = 0$.
    \item $\beta_1$: the slope — average increase in $Y$ associated with a one-unit increase in $X$.
\end{itemize}
\end{tcolorbox}

The least squares line is an \textbf{estimate} of the population regression line, just as the sample mean $\hat{\mu}$ estimates the population mean $\bar{\mu}$. The least squares line differs from the true line but is close, in the same way that the sample mean is different from but close to the population mean.

\subsection{Bias and Standard Error}

An estimator is \textbf{unbiased} if it does not systematically over- or under-estimate the true parameter. The sample mean $\hat{\mu}$ is an unbiased estimator of $\bar{\mu}$: over many datasets, the average of $\hat{\mu}$ coincides with $\bar{\mu}$.

To quantify how far a single estimator $\hat{\mu}$ will be from $\bar{\mu}$, we use the \textbf{Standard Error}:

\begin{tcolorbox}[colback=blue!5!white,colframe=blue!75!black,title=\textbf{Standard Error}]
\[
    \mathrm{Var}(\hat{\mu}) = \mathrm{SE}(\hat{\mu})^2 = \frac{\sigma^2}{n}
\]
where $\sigma$ is the standard deviation of each $y_i$ in $Y$. The more observations we have, the smaller the standard error.
\end{tcolorbox}

Analogously, the standard errors of $\hat{\beta}_0$ and $\hat{\beta}_1$ are:
\[
    \mathrm{SE}(\hat{\beta}_0)^2 = \sigma^2\!\left[\frac{1}{n} + \frac{\bar{x}^2}{\sum_{i=1}^{n}(x_i - \bar{x})^2}\right], \qquad
    \mathrm{SE}(\hat{\beta}_1)^2 = \frac{\sigma^2}{\sum_{i=1}^{n}(x_i - \bar{x})^2}
\]
where $\sigma^2 = \mathrm{Var}(\epsilon)$. Since $\sigma^2$ is not known, it is estimated from data via the \textbf{Residual Standard Error}:
\[
    \mathrm{RSE} = \sqrt{\frac{\mathrm{RSS}}{n-2}}
\]

\subsection{Confidence Intervals}

A \textbf{95\% confidence interval} is a range of values such that there is a 95\% probability that the range contains the true unknown value of the parameter. For $\beta_i$, assuming a Gaussian distribution, the 95\% CI approximately takes the form:
\[
    \hat{\beta}_i \pm 2 \cdot \mathrm{SE}(\hat{\beta}_i)
\]
That is, there is approximately a 95\% chance that the interval $\bigl[\hat{\beta}_i - 2\,\mathrm{SE}(\hat{\beta}_i),\; \hat{\beta}_i + 2\,\mathrm{SE}(\hat{\beta}_i)\bigr]$ contains the true value of $\beta_i$.

\subsection{Hypothesis Testing}

Standard errors can also be used to perform \textbf{hypothesis tests} on the coefficients. The most common test is:
\[
    H_0: \beta_1 = 0 \quad \text{(no relationship between } X \text{ and } Y\text{)}
\]
versus
\[
    H_a: \beta_1 \neq 0 \quad \text{(some relationship exists)}
\]
To test $H_0$, we need to assess whether $\hat{\beta}_1$ is sufficiently far from zero. We compute the \textbf{$t$-statistic}:
\[
    t = \frac{\hat{\beta}_1 - 0}{\mathrm{SE}(\hat{\beta}_1)}
\]
which measures the number of standard deviations that $\hat{\beta}_1$ is away from 0. We then compute the \textbf{$p$-value}: the probability of observing any value of the $t$-statistic equal to $|t|$ or larger, assuming $\beta_1 = 0$. A \textbf{small $p$-value} implies there is an association between the predictor and the response, so we can \textbf{reject the null hypothesis}. Typical cut-offs are $0.01$ and $0.05$.

\subsection{Assessing Model Accuracy: RSE and \texorpdfstring{$R^2$}{R2}}

Two standard measures quantify how well the model fits the data.

\begin{tcolorbox}[colback=blue!5!white,colframe=blue!75!black,title=\textbf{Residual Standard Error (RSE)}]
\[
    \mathrm{RSE} = \sqrt{\frac{1}{n-2}\,\mathrm{RSS}} = \sqrt{\frac{1}{n-2}\sum_{i=1}^{n}(y_i - \hat{y}_i)^2}
\]
RSE estimates the standard deviation of $\epsilon$. It measures the \textbf{lack of fit} of the linear model to the data, in the units of $Y$.
\end{tcolorbox}

\begin{tcolorbox}[colback=blue!5!white,colframe=blue!75!black,title=\textbf{$R^2$ Statistic}]
\[
    R^2 = \frac{\mathrm{TSS} - \mathrm{RSS}}{\mathrm{TSS}} = 1 - \frac{\mathrm{RSS}}{\mathrm{TSS}}
\]
where $\mathrm{TSS} = \sum(y_i - \bar{y})^2$ is the \textbf{total sum of squares}. $R^2$ is a proportion (between 0 and 1), independent of the scale of $Y$.
\end{tcolorbox}

Interpretation:
\begin{itemize}
    \item TSS measures total variance in $Y$ \textit{before} regression; RSS measures the variance remaining \textit{after} regression.
    \item $R^2$ measures the \textbf{proportion of variability in $Y$ explained by $X$}.
    \item $R^2 \approx 1$: the regression explains most of the variability — good linear fit.
    \item $R^2 \approx 0$: the regression explains little — the linear model is wrong or $\sigma^2$ is high.
\end{itemize}

In simple linear regression, $R^2 = r^2$ where $r = \mathrm{Cor}(X, Y)$ is the sample correlation:
\[
    \mathrm{Cor}(X, Y) = \frac{\sum_{i=1}^{n}(x_i - \bar{x})(y_i - \bar{y})}{\sqrt{\sum_{i=1}^{n}(x_i - \bar{x})^2}\sqrt{\sum_{i=1}^{n}(y_i - \bar{y})^2}}
\]


\section{Multiple Linear Regression}

\subsection{The Model}

With $p$ predictors, the model generalizes to:
\[
    Y = \beta_0 + \beta_1 X_1 + \cdots + \beta_p X_p + \epsilon
\]
$\beta_j$ represents the \textbf{average effect on $Y$ of a one-unit increase in $X_j$, holding all other predictors fixed}. For the advertising example:
\[
    \text{sales} = \beta_0 + \beta_1 \times \text{TV} + \beta_2 \times \text{radio} + \beta_3 \times \text{newspaper} + \epsilon
\]

\subsection{Estimating the Coefficients}

$\beta_0, \beta_1, \ldots, \beta_p$ are unknown and must be estimated. Given estimates $\hat{\beta}_0, \hat{\beta}_1, \ldots, \hat{\beta}_p$, predictions are made via:
\[
    \hat{y} = \hat{f}(X_p) = \hat{\beta}_0 + \sum_{j=1}^{p} \hat{\beta}_j x_j
\]
We choose $\hat{\beta}_0, \hat{\beta}_1, \ldots, \hat{\beta}_p$ to minimize the RSS:
\[
    \mathrm{RSS} = \sum_{i=1}^{n}(y_i - \hat{y}_i)^2 = \sum_{i=1}^{n}\!\left(y_i - \hat{\beta}_0 - \sum_{j=1}^{p}\hat{\beta}_j x_{ij}\right)^{\!2}
\]
Unlike simple linear regression, the minimizers do not have a simple closed form and are computed via matrix algebra (normal equations).

\begin{remark}
Correlation between predictors can make individual coefficient estimates difficult to interpret. In the advertising example, newspaper has a significant positive coefficient in simple regression ($p < 0.001$) but a near-zero, non-significant coefficient ($p = 0.86$) in multiple regression. This is because newspaper is correlated with radio ($r = 0.354$): in markets with high radio budgets, newspaper budgets also tend to be high, and sales are driven by radio, not newspaper.
\end{remark}

\begin{tcolorbox}[colback=orange!5!white,colframe=orange!75!black]
\textbf{Correlation $\neq$ causality.} A predictor may appear significant in simple regression only because it is correlated with the true driving variable. Multiple regression helps disentangle these effects but does not establish causation.
\end{tcolorbox}

\subsection{Important Questions in Multiple Regression}

\subsubsection{Question 1: Is at least one predictor useful?}

We test the joint null hypothesis $H_0: \beta_1 = \beta_2 = \cdots = \beta_p = 0$ against $H_a$: at least one $\beta_j \neq 0$ using the \textbf{$F$-statistic}:
\[
    F = \frac{(\mathrm{TSS} - \mathrm{RSS})/p}{\mathrm{RSS}/(n-p-1)}
\]
where $\mathrm{TSS} = \sum(y_i - \bar{y})^2$ and $\mathrm{RSS} = \sum(y_i - \hat{y}_i)^2$. Small values of $F$ indicate no response-predictor relationship.

\subsubsection{Question 2: Which predictors are important?}

When $p$ is large, testing all $2^p$ subsets is infeasible. Three classical automated approaches exist:

\begin{itemize}
    \item \textbf{Forward selection}: start with the null model (intercept only). Fit $p$ simple regressions and add the variable with the lowest RSS. Continue adding variables that most reduce RSS, until a stopping rule is satisfied.
    \item \textbf{Backward selection}: start with all $p$ variables. Remove the variable with the largest $p$-value (least statistically significant). Refit with $p-1$ variables and repeat, until all remaining variables have $p$-value below a threshold.
    \item \textbf{Mixed selection}: start with no variables and add the best fitting variable one at a time. If the $p$-value for any current variable rises above a threshold as new variables are added, remove it. Continue until all variables in the model have a sufficiently low $p$-value.
\end{itemize}

\subsubsection{Question 3: Model fit}

In multiple regression, the RSE is defined as:
\[
    \mathrm{RSE} = \sqrt{\frac{1}{n-p-1}\,\mathrm{RSS}}
\]
Adding more predictors to the model will always increase $R^2$ (even if the additional predictor is weakly related to the response). Graphical summaries — such as residual plots — can reveal problems not visible from numerical statistics alone. For example, a non-linear pattern in the residuals suggests a \textbf{synergy} or \textbf{interaction effect} between predictors.

\subsubsection{Question 4: Predictions}

Three sources of uncertainty are associated with regression predictions:
\begin{enumerate}
    \item The least squares plane is only an estimate for the true population regression plane (different datasets yield different coefficients).
    \item Model bias: assuming a linear $f(X)$ is almost always an approximation of reality (reducible error).
    \item Even knowing the true $f(X)$, an irreducible error $\epsilon$ is always present — we use \textbf{prediction intervals} to account for this.
\end{enumerate}


\section{Considerations about the Regression Model}

\subsection{Qualitative Predictors}

\subsubsection{Two-Level Qualitative Predictors}

When a qualitative predictor (factor) has two levels, we create an \textbf{indicator} or \textbf{dummy variable} that takes two numerical values. For example, to encode gender in a credit card balance study:
\[
    x_i = \begin{cases} 1 & \text{if the } i\text{-th person is female} \\ 0 & \text{if the } i\text{-th person is male} \end{cases}
\]
The regression equation becomes:
\[
    y_i = \beta_0 + \beta_1 x_i + \epsilon_i = \begin{cases} \beta_0 + \beta_1 + \epsilon_i & \text{if female} \\ \beta_0 + \epsilon_i & \text{if male} \end{cases}
\]
Interpretation: $\beta_0$ is the average balance among males; $\beta_0 + \beta_1$ is the average among females; $\beta_1$ is the average difference in balance between females and males.

\subsubsection{Qualitative Predictors with More than Two Levels}

A single dummy variable cannot encode more than two levels. For a qualitative predictor with $K$ levels we create $K-1$ dummy variables. For instance, for a three-level ethnicity variable (Asian, Caucasian, African American):
\[
    x_{i1} = \begin{cases} 1 & \text{Asian} \\ 0 & \text{otherwise} \end{cases}
    \qquad
    x_{i2} = \begin{cases} 1 & \text{Caucasian} \\ 0 & \text{otherwise} \end{cases}
\]
The model becomes $y_i = \beta_0 + \beta_1 x_{i1} + \beta_2 x_{i2} + \epsilon_i$, which gives:
\[
    y_i = \begin{cases} \beta_0 + \beta_1 + \epsilon_i & \text{Asian} \\ \beta_0 + \beta_2 + \epsilon_i & \text{Caucasian} \\ \beta_0 + \epsilon_i & \text{African American} \end{cases}
\]
The level with no dummy variable is the \textbf{baseline}.

\subsection{Extensions of the Linear Model}

The standard linear model assumes the relationship between predictors and response is \textbf{linear and additive}:
\begin{itemize}
    \item \textbf{Additivity}: the effect of $X_j$ on $Y$ is independent of the other predictors.
    \item \textbf{Linearity}: the change in $Y$ due to a one-unit change in $X_j$ is constant.
\end{itemize}

\subsubsection{Interaction Effects}

The \textbf{synergy} (or \textbf{interaction}) effect occurs when the effect of one predictor depends on the value of another. To accommodate it, we include an \textbf{interaction term} $X_1 X_2$:
\[
    Y = \beta_0 + \beta_1 X_1 + \beta_2 X_2 + \beta_3 X_1 X_2 + \epsilon
\]
Defining $\tilde{\beta}_1 = (\beta_1 + \beta_3 X_2)$, this can be rewritten as $Y = \beta_0 + \tilde{\beta}_1 X_1 + \beta_2 X_2 + \epsilon$. The effect of $X_1$ on $Y$ is no longer constant: adjusting $X_2$ changes the impact of $X_1$ on $Y$.

\subsubsection{Non-linear Relationships}

\textbf{Polynomial regression} extends the linear model to accommodate non-linear relationships while remaining a linear model in the coefficients:
\[
    \text{mpg} = \beta_0 + \beta_1 \times \text{horsepower} + \beta_2 \times \text{horsepower}^2 + \epsilon
\]
The second predictor is $X_2 = \text{horsepower}^2$; the model is still linear in $\beta_0, \beta_1, \beta_2$.

\subsection{Potential Problems}

\subsubsection{1. Non-linearity of the Data}

The linear model assumes a linear relationship between $X$ and $Y$. \textbf{Residual plots} (plotting $e_i = y_i - \hat{y}_i$ versus fitted values $\hat{y}_i$) help identify non-linearity: a pattern in the residuals signals that the linear assumption is violated. A simple remedy is to apply \textbf{non-linear transformations} of the predictors, such as $\log X$ or $\sqrt{X}$.

\subsubsection{2. Correlation of Error Terms}

An important assumption is that the error terms $\epsilon_1, \epsilon_2, \ldots, \epsilon_n$ are \textbf{uncorrelated}. If there is correlation among errors:
\begin{itemize}
    \item Estimated standard errors will underestimate the true standard error.
    \item Confidence intervals will be narrower than they should be.
    \item $p$-values will be lower than they should be — parameters may be declared erroneously significant.
\end{itemize}
Correlated errors create an \textbf{unwarranted sense of confidence} in the model. This problem frequently arises in time series data.

\subsubsection{3. Non-constant Variance of Error Terms}

Linear regression also assumes \textbf{homoscedasticity}: $\mathrm{Var}(\epsilon_i) = \sigma^2$ (constant). When the variance of the errors increases with the fitted values, this is called \textbf{heteroscedasticity}. It is identifiable from a funnel-shaped residual plot. A common solution is to transform the response with a concave function such as $\log(Y)$.

\subsubsection{4. Outliers}

An \textbf{outlier} is a point $y_i$ that is far from the value predicted by the model (arising, for instance, from incorrect recording). Outliers have little effect on the least squares fit itself but can cause problems with \textbf{confidence intervals and $p$-values}: for example, an outlier may inflate the RSE, making all confidence intervals wider and $p$-values less significant.

\subsubsection{5. Collinearity}

\begin{tcolorbox}[colback=blue!5!white,colframe=blue!75!black,title=\textbf{Collinearity}]
\textbf{Collinearity} refers to the situation in which two or more predictor variables are closely related to each other.
\end{tcolorbox}

Collinearity makes it difficult to separate out the individual effects of collinear variables on the response. Geometrically, the RSS contour becomes elongated: there are many pairs $(\hat{\beta}_{\text{limit}}, \hat{\beta}_{\text{rating}})$ with a similar RSS, so the minimum is not well-determined. Consequences:
\begin{itemize}
    \item Estimates $\hat{\beta}_j$ become very uncertain (high standard errors).
    \item The $t$-statistic for each collinear predictor declines, inflating $p$-values — a truly associated variable may fail to have a small $p$-value.
\end{itemize}
Collinearity can be detected from the \textbf{correlation matrix} of predictors. When more than two variables are involved, the \textbf{Variance Inflation Factor (VIF)} is used; its presence motivates variable selection methods such as those covered in subsequent chapters.


% ========================================
\part{Cross-Validation}
% ========================================

\section{Resampling Methods}

\textbf{Resampling methods} involve repeatedly drawing samples from a training set and refitting a model of interest on each sample in order to obtain additional information about that model. For example, we can estimate the variability of a linear regression fit by drawing different samples from the training data and fitting a linear regression to each one. The main drawback is computational cost — fitting the same statistical method multiple times using different subsets of the training data — but this is generally not prohibitive.

The two most commonly used resampling methods are \textbf{cross-validation} and \textbf{bootstrapping}. Cross-validation is used to estimate the test error associated with a given statistical learning method and to select the appropriate level of model flexibility.

\section{Test Error vs.\ Training Error}

\begin{itemize}
    \item The \textbf{test error} is the average error that results from using a statistical learning method to predict the response on a \textbf{new observation} — a measurement not used for training. A method can be considered warranted if it produces models with low test error.
    \item The \textbf{training error} can be easily calculated by applying the method to the observations used in its training. Unfortunately, the training error rate often \textbf{underestimates} the test error rate considerably.
\end{itemize}

When a dedicated test set is not available, we consider methods that estimate the test error by \textbf{holding out} a subset of training observations from the fitting process and then evaluating the error on that held-out subset.

\section{The Validation Set Approach}

\begin{tcolorbox}[colback=blue!5!white,colframe=blue!75!black,title=\textbf{Validation Set Approach}]
Randomly divide the available observations into two parts: a \textbf{training set} and a \textbf{validation set} (or hold-out set) of approximately equal size. Fit the model on the training set, then evaluate predictions on the validation set. The resulting \textbf{validation set error rate} — typically the MSE — provides an estimate of the test error rate.
\end{tcolorbox}

\subsection{Drawbacks}

\begin{enumerate}
    \item The validation estimate of the test error rate can be \textbf{highly variable}, depending on which observations happen to be included in the training set and which in the validation set.
    \item Only a subset of observations is used to fit the model. Since statistical methods generally perform better with more data, the \textbf{validation set error rate may overestimate} the test error rate for the model fit on the entire dataset.
\end{enumerate}

\textit{Example}: on the Auto dataset, using the validation set approach to estimate the test MSE for a polynomial regression of \texttt{mpg} on \texttt{horsepower}, the estimated MSE varies substantially across different random splits of the data (right panel of the plot shows the high variability across ten random splits).

\section{Leave-One-Out Cross-Validation (LOOCV)}

\begin{tcolorbox}[colback=blue!5!white,colframe=blue!75!black,title=\textbf{Leave-One-Out Cross-Validation}]
LOOCV involves splitting the observations into two parts. Instead of two comparable-size subsets, a \textbf{single observation} $(x_1, y_1)$ is used as the validation set, and the remaining $n-1$ observations $(x_2, y_2), \ldots, (x_n, y_n)$ form the training set.
\end{tcolorbox}

The model is fit on $n-1$ observations and a prediction $\hat{y}_1$ is made for the excluded observation using $x_1$. Since $(x_1, y_1)$ was not used in fitting:
\[
    \mathrm{MSE}_1 = (y_1 - \hat{y}_1)^2
\]
is an approximately unbiased estimate of the test error. However, it is a \textbf{poor estimate} on its own because it is based on a single observation and hence highly variable.

Repeating the procedure by selecting each $(x_i, y_i)$ in turn as the validation observation produces $n$ squared errors $\mathrm{MSE}_1, \ldots, \mathrm{MSE}_n$. The LOOCV estimate for the test MSE averages these:
\[
    \mathrm{CV}(n) = \frac{1}{n} \sum_{i=1}^{n} \mathrm{MSE}_i
\]

\subsection{Advantages over the Validation Set Approach}

\begin{enumerate}
    \item \textbf{Far less bias}: each training set contains $n-1$ observations, almost as many as the full dataset. LOOCV therefore tends not to overestimate the test error rate as much as the validation set approach.
    \item \textbf{No randomness}: LOOCV always yields the same result — there is no randomness in the training/validation split.
\end{enumerate}

LOOCV is a very general method and can be used with any kind of predictive modeling (logistic regression, linear discriminant analysis, etc.).

\section{\texorpdfstring{$k$}{k}-Fold Cross-Validation}

\begin{tcolorbox}[colback=blue!5!white,colframe=blue!75!black,title=\textbf{$k$-Fold Cross-Validation}]
Randomly divide the set of observations into $k$ groups (folds) of approximately equal size.
\begin{enumerate}
    \item Treat the first fold as the validation set and fit the method on the remaining $k-1$ folds; compute $\mathrm{MSE}_1$ on the validation fold.
    \item Repeat $k$ times, each time using a different fold as the validation set.
    \item This yields $k$ estimates $\mathrm{MSE}_1, \ldots, \mathrm{MSE}_k$. The $k$-fold CV estimate is:
\end{enumerate}
\[
    \mathrm{CV}(k) = \frac{1}{k} \sum_{i=1}^{k} \mathrm{MSE}_i
\]
\end{tcolorbox}

LOOCV is a special case of $k$-fold CV in which $k = n$.

\subsection{Advantages of \texorpdfstring{$k$}{k}-fold CV over LOOCV}

In practice, $k$-fold CV is typically performed with $k = 5$ or $k = 10$. The advantages over LOOCV are:

\begin{enumerate}
    \item \textbf{Computational}: LOOCV requires fitting the model $n$ times, which can be prohibitively expensive for large $n$. $k$-fold CV with $k = 10$ requires fitting only 10 times.
    \item \textbf{Bias-variance trade-off}: each random split of observations estimates the test error differently. The variability in $k$-fold CV estimates is lower than the variability in the validation set approach, though higher than LOOCV's variability.
\end{enumerate}

\subsection{Bias-Variance Trade-off in Cross-Validation}

The choice of $k$ in $k$-fold CV involves a bias-variance trade-off:

\begin{itemize}
    \item \textbf{Validation set approach}: each training set contains only about half the observations $\Rightarrow$ \textbf{high bias} (overestimates test error), low variance issues come from using a single split.
    \item \textbf{LOOCV}: each training set contains $n-1$ observations $\Rightarrow$ \textbf{approximately unbiased} estimates of the test error. However, LOOCV has \textbf{higher variance} than $k$-fold CV.
    \item \textbf{$k$-fold CV} ($k < n$): each training set contains $(k-1)n/k$ observations $\Rightarrow$ \textbf{intermediate level of bias}.
\end{itemize}

\subsubsection{Why does LOOCV have higher variance?}

When performing LOOCV we average the outputs of $n$ fitted models, each trained on an almost identical set of observations. These outputs are therefore \textbf{highly (positively) correlated} with each other. In contrast, with $k$-fold CV ($k < n$) we average the outputs of $k$ fitted models whose training sets overlap less — they are \textbf{less correlated}. Since the mean of many highly correlated quantities has higher variance than the mean of less correlated quantities, the test error estimate from LOOCV tends to have higher variance than that from $k$-fold CV.

\begin{tcolorbox}[colback=green!5!white,colframe=green!75!black]
From the perspective of \textbf{bias reduction}, LOOCV should be preferred to $k$-fold CV. However, accounting for \textbf{variance}, $k$-fold CV with $k = 5$ or $k = 10$ is typically the preferred choice: these values have been shown empirically to yield test error rate estimates that suffer neither from excessively high bias nor from very high variance.
\end{tcolorbox}

\subsection{CV Curves and the True Test MSE}

When examining real data, the true test MSE is unknown, so it is difficult to assess the accuracy of the CV estimate directly. On simulated data, where the true test MSE can be computed, CV curves have the \textbf{correct general shape} — they identify the right level of flexibility — but they tend to \textbf{underestimate} the true test MSE slightly. For model selection purposes (choosing among competing models or selecting the degree of flexibility), what matters is the shape of the curve rather than its absolute value.


% ========================================
\part{Linear Model Selection}
% ========================================

\section{Introduction}

\subsection{Why Use Alternative Fitting Procedures?}

The standard linear model $Y = \beta_0 + \sum_{j=1}^p \beta_j X_j + \epsilon$ fit by least squares is a workhorse of statistical learning: it offers strong inference guarantees and, on real-world problems, is often surprisingly competitive with non-linear methods. However, plain least squares can be improved by replacing it with \textbf{alternative fitting procedures}, for two reasons.

\subsubsection{Prediction Accuracy}

\begin{itemize}
    \item If $n \gg p$ (observations much larger than variables), least squares estimates have low variance and perform well on test observations.
    \item If $n$ is \textbf{not much larger than} $p$, there can be substantial variability in the least squares fit, leading to overfitting and poor predictions on unseen data.
    \item If $p > n$, there is \textbf{no unique least squares estimate}: variance is infinite and the method cannot be used at all.
\end{itemize}

\subsubsection{Model Interpretability}

In multiple regression, some predictors are often \textbf{not associated with the response}, introducing unnecessary complexity. Setting the corresponding coefficient estimates to zero produces a more interpretable model. This is the goal of \textbf{feature selection} (or \textbf{variable selection}).

\subsection{The Prediction Accuracy vs.\ Interpretability Trade-off}

\begin{center}
\begin{tabular}{lll}
 & \textbf{Fewer predictors} & \textbf{More predictors} \\
\hline
Interpretability & Easy & Difficult \\
Overfitting risk & Low & High \\
Prediction accuracy & Low & High \\
Computational cost & Low & High \\
\end{tabular}
\end{center}

\subsection{Three Classes of Approaches}

There are three main strategies for fitting $Y = \beta_0 + \sum_{i=1}^p \beta_i X_i + \epsilon$ beyond plain least squares:

\begin{itemize}
    \item \textbf{Subset Selection}: select a subset of the $p$ predictors and fit the model by least squares on the reduced set.
    \item \textbf{Shrinkage} (regularization): fit the model on all $p$ predictors but shrink the estimated coefficients towards zero. This reduces variance; some coefficients may be shrunk to exactly zero, performing variable selection implicitly.
    \item \textbf{Dimension Reduction}: project the $p$ predictors into an $M$-dimensional subspace ($M < p$) by computing $M$ linear combinations of the variables, then fit a linear regression model on these $M$ projections using least squares.
\end{itemize}


\section{Best Subset Selection}

\subsection{Algorithm}

Best subset selection fits a separate least squares regression for \textbf{every possible combination} of the $p$ predictors.

\begin{enumerate}
    \item Let $\mathcal{M}_0$ denote the \textbf{null model} — no predictors, predicts the sample mean for every observation.
    \item For $k = 1, \ldots, p$:
        \begin{enumerate}[label=(\alph*)]
            \item Fit all $\binom{p}{k} = \frac{p!}{k!(p-k)!}$ models that contain exactly $k$ predictors.
            \item Choose the best among these models (smallest RSS / highest $R^2$) and call it $\mathcal{M}_k$.
        \end{enumerate}
    \item Select a single best model from $\mathcal{M}_0, \ldots, \mathcal{M}_p$ using \textbf{cross-validated prediction error}, $C_p$, AIC, BIC, or adjusted $R^2$.
\end{enumerate}

\begin{remark}
RSS decreases monotonically and $R^2$ increases monotonically as $k$ grows — they cannot be used to select among $\mathcal{M}_0, \ldots, \mathcal{M}_p$ directly, as they always favour the largest model. Step 3 requires a criterion that accounts for model complexity.
\end{remark}

\subsection{Computational Limitation}

The total number of models is $2^p$. This grows rapidly with $p$:
\begin{itemize}
    \item $p = 10$: $\approx 1{,}000$ models.
    \item $p = 20$: $> 1{,}000{,}000$ models.
\end{itemize}
Best subset selection is \textbf{computationally infeasible for large $p$}.


\section{Forward Stepwise Selection}

\subsection{Algorithm}

Forward stepwise selection is a \textbf{computationally efficient alternative} that considers a much smaller set of models. It starts with no predictors and adds them one at a time.

\begin{enumerate}
    \item Let $\mathcal{M}_0$ denote the \textbf{null model} with no predictors.
    \item For $k = 0, \ldots, p-1$:
        \begin{enumerate}[label=(\alph*)]
            \item Consider all $p - k$ models that augment $\mathcal{M}_k$ with one additional predictor.
            \item Choose the best among these (smallest RSS or highest $R^2$) and call it $\mathcal{M}_{k+1}$.
        \end{enumerate}
    \item Select a single best model from $\mathcal{M}_0, \ldots, \mathcal{M}_p$ using cross-validation, $C_p$, AIC, BIC, or adjusted $R^2$.
\end{enumerate}

\subsection{Computational Cost}

Forward stepwise selection fits $1 + \sum_{k=0}^{p-1}(p - k) = 1 + p(p+1)/2$ models in total. For $p = 20$: only \textbf{211 models} versus 1,048,576 for best subset selection.

\begin{remark}
Forward stepwise selection does \textbf{not} guarantee finding the globally best model among all $2^p$ subsets, since it follows a greedy path. However, it can be applied even when $n < p$, whereas best subset selection cannot.
\end{remark}


\section{Backward Stepwise Selection}

\subsection{Algorithm}

Backward stepwise selection is an efficient alternative that begins with the \textbf{full model} and iteratively removes the least useful predictor.

\begin{enumerate}
    \item Let $\mathcal{M}_p$ denote the \textbf{full model} with all $p$ predictors.
    \item For $k = p, p-1, \ldots, 1$:
        \begin{enumerate}[label=(\alph*)]
            \item Consider all $k$ models that contain all but one of the predictors in $\mathcal{M}_k$ (i.e., $k-1$ predictors each).
            \item Choose the best (smallest RSS) and call it $\mathcal{M}_{k-1}$.
        \end{enumerate}
    \item Select a single best model from $\mathcal{M}_0, \ldots, \mathcal{M}_p$ using cross-validation, $C_p$, AIC, BIC, or adjusted $R^2$.
\end{enumerate}

Like forward stepwise, backward stepwise searches through only $1 + p(p+1)/2$ models. It \textbf{requires $n > p$} so that the full model can be fit. Forward stepwise can be used when $n < p$ and is therefore the only viable stepwise method in the high-dimensional setting.

\subsection{Comparison of the Three Algorithms}

The three algorithms generally give \textbf{similar but not identical} models. A \textbf{hybrid approach} also exists: variables are added sequentially (like forward selection) but any variable whose $p$-value rises above a threshold after new variables are added is removed (like backward selection), attempting to mimic best subset selection while retaining the computational advantages of stepwise methods.


\section{Choosing the Optimal Model}

All three subset selection procedures produce a sequence of models $\mathcal{M}_0, \ldots, \mathcal{M}_p$. To select the best among them we need to estimate the \textbf{test error}. RSS and $R^2$ track training error and cannot be used directly. There are two approaches.

\subsection{Approach 1: Adjusting the Training Error}

The training MSE ($= \mathrm{RSS}/n$) systematically underestimates the test MSE. Four criteria adjust for model size $d$ (number of predictors).

\subsubsection{Mallow's \texorpdfstring{$C_p$}{Cp}}

\[
    C_p = \frac{1}{n}\bigl(\mathrm{RSS} + 2d\hat{\sigma}^2\bigr)
\]
where $\hat{\sigma}^2$ is an estimate of $\mathrm{Var}(\epsilon)$. The penalty $2d\hat{\sigma}^2$ increases with $d$, correcting for the tendency of training RSS to underestimate test error. We select the model with the \textbf{lowest} $C_p$.

\subsubsection{AIC — Akaike Information Criterion}

Defined for models fit by maximum likelihood. For a multivariate linear model with Gaussian errors, maximum likelihood and least squares are equivalent, giving:
\[
    \mathrm{AIC} = \frac{1}{n\hat{\sigma}^2}\bigl(\mathrm{RSS} + 2d\hat{\sigma}^2\bigr)
\]
AIC and $C_p$ are proportional in this setting and select the same model.

\subsubsection{BIC — Bayesian Information Criterion}

\[
    \mathrm{BIC} = \frac{1}{n}\bigl(\mathrm{RSS} + \log(n)\,d\hat{\sigma}^2\bigr)
\]
BIC replaces the factor $2$ in $C_p$ with $\log(n)$. Since $\log(n) > 2$ for $n > 7$, BIC places a \textbf{heavier penalty} on models with many variables and tends to select \textbf{smaller models} than $C_p$. We select the model with the lowest BIC.

\subsubsection{Adjusted \texorpdfstring{$R^2$}{R2}}

Standard $R^2$ always increases as more variables are added. The adjusted $R^2$ corrects for this:
\[
    \text{Adjusted } R^2 = 1 - \frac{\mathrm{RSS}/(n - d - 1)}{\mathrm{TSS}/(n-1)}
\]
A large adjusted $R^2$ indicates a small test error. Unlike $R^2$, adding an unnecessary variable will decrease adjusted $R^2$ because the decrease in RSS is small relative to the increase in $d$. We select the model with the \textbf{highest} adjusted $R^2$.

\subsection{Approach 2: Direct Estimation of Test Error}

We can directly estimate the test error using the \textbf{validation set approach} or \textbf{cross-validation}. Plotting the validation set error or $k$-fold CV error as a function of the number of predictors $d$ typically shows a clear minimum, identifying the optimal model size.

\begin{tcolorbox}[colback=green!5!white,colframe=green!75!black]
All four indirect criteria ($C_p$, AIC, BIC, adjusted $R^2$) and the direct CV approach tend to \textbf{agree} in practice, identifying the same optimal model. When they disagree, direct CV is generally preferred as it makes fewer assumptions about the error distribution.
\end{tcolorbox}


%=============================================================
\part{Regularization}
%=============================================================

Subset selection methods (best subset, forward/backward stepwise) reduce model complexity by selecting a \textbf{subset of predictors}. An alternative strategy is to fit a model containing \textbf{all $p$ predictors} but to \textbf{constrain} (regularize) the coefficient estimates, shrinking them toward zero. This can significantly \textbf{reduce variance} at the cost of a small increase in bias, often leading to better prediction accuracy. The two best-known shrinkage techniques are \textbf{ridge regression} and the \textbf{lasso}.

%-------------------------------------------------------------
\section{Shrinkage Methods}
%-------------------------------------------------------------

\begin{tcolorbox}[colback=blue!5!white,colframe=blue!75!black,title=\textbf{Shrinkage Methods}]
Shrinkage methods fit a model that contains all $p$ predictors using a technique that \textbf{constrains} or \textbf{regularizes} the coefficient estimates, equivalently \textbf{shrinking} them toward zero. Shrinking the coefficient estimates can significantly \textbf{reduce their variance}.
\end{tcolorbox}

%-------------------------------------------------------------
\section{Ridge Regression}
%-------------------------------------------------------------

\subsection{Formulation}

Recall that the ordinary least squares (OLS) estimates $\hat{\beta}_0, \hat{\beta}_1, \ldots, \hat{\beta}_p$ minimize the residual sum of squares:
\[
    \mathrm{RSS} = \sum_{i=1}^{n} \left(y_i - \beta_0 - \sum_{j=1}^{p} \beta_j x_{ij}\right)^{\!2}
\]

\textbf{Ridge regression} is very similar, but the coefficient estimates $\hat{\beta}^R_\lambda$ minimize a \textbf{penalized} objective:

\begin{tcolorbox}[colback=blue!5!white,colframe=blue!75!black,title=\textbf{Ridge Regression Objective}]
\[
    \sum_{i=1}^{n} \left(y_i - \beta_0 - \sum_{j=1}^{p} \beta_j x_{ij}\right)^{\!2} + \lambda \sum_{j=1}^{p} \beta_j^2
    \;=\; \mathrm{RSS} + \lambda \sum_{j=1}^{p} \beta_j^2
\]
where $\lambda \geq 0$ is a \textbf{tuning parameter}.
\end{tcolorbox}

Like OLS, ridge regression seeks coefficients that fit the data well by making the RSS small. However, the second term --- the \textbf{shrinkage penalty} $\lambda \sum_j \beta_j^2$ (an $\ell_2$ penalty) --- is small when $\beta_1, \ldots, \beta_p$ are close to zero.

\subsection{The Tuning Parameter $\lambda$}

The parameter $\lambda$ controls the \textbf{relative impact} of the two terms:
\begin{itemize}
    \item $\lambda = 0$: the penalty has no effect $\Rightarrow$ ridge regression coincides with \textbf{least squares}.
    \item $\lambda \to \infty$: the penalty dominates $\Rightarrow$ all coefficient estimates are driven toward zero.
\end{itemize}
Ridge regression produces a \textbf{different set of coefficient estimates} $\hat{\beta}^R_\lambda$ for each value of $\lambda$. As $\lambda$ increases from zero, the ridge coefficient estimates shrink toward zero.

\begin{remark}
The shrinkage penalty is \textbf{not} applied to $\beta_0$, which is simply the mean of the response when all predictors are centered. Only $\beta_1, \ldots, \beta_p$ are penalized.
\end{remark}

\subsection{Scale Sensitivity and Standardization}

Unlike least squares, ridge regression is \textbf{not scale invariant}: multiplying a predictor $X_j$ by a constant $c$ changes the ridge coefficient estimates because the penalty $\lambda \sum_j \beta_j^2$ depends on the scale of each predictor. Therefore, one should \textbf{standardize the predictors} before applying ridge regression:
\[
    \tilde{x}_{ij} = \frac{x_{ij} - \bar{x}_j}{\sigma_j}
    \qquad \text{where} \qquad
    \sigma_j = \sqrt{\frac{1}{n}\sum_{i=1}^{n}(x_{ij} - \bar{x}_j)^2}
\]
so that all predictors are on the same scale and the final fit does not depend on the units of measurement.

\subsection{The $\ell_2$ Norm Perspective}

Let the $\ell_2$ norm of the coefficient vector be $\|\beta\|_2 = \sqrt{\sum_{j=1}^{p} \beta_j^2}$. The ratio $\|\hat{\beta}^R_\lambda\|_2 / \|\hat{\beta}\|_2$ ranges from 1 (when $\lambda = 0$, recovering OLS) to 0 (when $\lambda = \infty$). This ratio measures the \textbf{amount by which} the ridge regression coefficients have been shrunk toward zero relative to the OLS solution.

\subsection{Bias-Variance Trade-off in Ridge Regression}

The advantage of ridge regression over least squares is rooted in the \textbf{bias-variance trade-off}:
\begin{itemize}
    \item As $\lambda$ increases, the \textbf{flexibility} of the ridge regression fit \textbf{decreases}: variance decreases but bias increases.
    \item For an intermediate value of $\lambda$, the MSE can be \textbf{considerably lower} than the OLS MSE, because the reduction in variance more than compensates for the increase in bias.
\end{itemize}

\subsection{When to Use Ridge Regression}

\begin{tcolorbox}[colback=green!5!white,colframe=green!75!black,title=\textbf{When Ridge Regression Helps}]
\begin{itemize}
    \item When the relationship between response and predictors is close to \textbf{linear}, the OLS estimates have low bias but may have \textbf{high variance} --- a small change in the training data can cause large changes in the coefficient estimates.
    \item In particular, when $p$ is almost as large as $n$, OLS estimates are \textbf{extremely variable}.
    \item If $p > n$, OLS does not even have a \textbf{unique solution}, whereas ridge regression performs well by trading a small increase in bias for a large decrease in variance.
    \item Ridge regression has substantial \textbf{computational advantages} over best subset selection: fitting $2^p$ models vs.\ a single model for any value of $\lambda$.
\end{itemize}
\end{tcolorbox}

\begin{remark}
With \textbf{multicollinearity} among predictors, the OLS loss surface is a long valley (a ``ridge''), making the solution unstable. The $\ell_2$ penalty turns this ridge into a \textbf{more regular} (better conditioned) function, stabilizing the estimates.
\end{remark}

%-------------------------------------------------------------
\section{The Lasso}
%-------------------------------------------------------------

\subsection{Motivation}

Ridge regression has one important \textbf{disadvantage}: it includes \textbf{all $p$ predictors} in the final model. The penalty shrinks all coefficients toward zero, but it does \textbf{not} set any of them \textbf{exactly to zero} (unless $\lambda = \infty$). This is not a problem for prediction accuracy, but it can hinder \textbf{model interpretability}, especially when $p$ is large.

\subsection{Formulation}

The \textbf{Lasso} (Least Absolute Shrinkage and Selection Operator) overcomes this disadvantage. The lasso coefficient estimates $\hat{\beta}^L_\lambda$ minimize:

\begin{tcolorbox}[colback=blue!5!white,colframe=blue!75!black,title=\textbf{Lasso Objective}]
\[
    \sum_{i=1}^{n} \left(y_i - \beta_0 - \sum_{j=1}^{p} \beta_j x_{ij}\right)^{\!2} + \lambda \sum_{j=1}^{p} |\beta_j|
    \;=\; \mathrm{RSS} + \lambda \sum_{j=1}^{p} |\beta_j|
\]
The lasso uses an $\ell_1$ penalty instead of the $\ell_2$ penalty: $\|\beta\|_1 = \sum_{j=1}^{p} |\beta_j|$.
\end{tcolorbox}

\subsection{Variable Selection}

The $\ell_1$ penalty has the effect of \textbf{forcing some coefficient estimates to be exactly zero} when $\lambda$ is sufficiently large. As a consequence, the lasso performs \textbf{variable selection}, yielding \textbf{sparse models} that involve only a subset of the predictors.

\begin{itemize}
    \item $\lambda = 0$: the lasso gives the \textbf{OLS} fit.
    \item $\lambda$ sufficiently large: the lasso gives the \textbf{null model} (all coefficients equal to zero).
    \item Between these extremes, the lasso produces models with \textbf{any number of variables}, depending on $\lambda$.
\end{itemize}

In contrast, ridge regression always includes all variables in the model (with magnitude depending on $\lambda$). Selecting a good value for the lasso parameter is critical --- as with ridge regression, \textbf{cross-validation} is used to select the optimal $\lambda$.

%-------------------------------------------------------------
\section{Constrained Formulation}
%-------------------------------------------------------------

The lasso and ridge regression objectives can be equivalently written as \textbf{constrained optimization} problems:

\begin{tcolorbox}[colback=blue!5!white,colframe=blue!75!black,title=\textbf{Constrained Formulations}]
\textbf{Lasso:}
\[
    \min_\beta \left\{\sum_{i=1}^{n} \left(y_i - \beta_0 - \sum_{j=1}^{p} \beta_j x_{ij}\right)^{\!2}\right\}
    \qquad \text{subject to} \qquad
    \sum_{j=1}^{p} |\beta_j| \leq s
\]
\textbf{Ridge:}
\[
    \min_\beta \left\{\sum_{i=1}^{n} \left(y_i - \beta_0 - \sum_{j=1}^{p} \beta_j x_{ij}\right)^{\!2}\right\}
    \qquad \text{subject to} \qquad
    \sum_{j=1}^{p} \beta_j^2 \leq s
\]
\end{tcolorbox}

We seek coefficient estimates that make the RSS as small as possible, subject to the requirement that the $\ell_1$ (lasso) or $\ell_2$ (ridge) norm does not exceed a \textbf{budget} $s$. The budget $s$ absorbs the role of $\lambda$: when $s$ is extremely large, the constraint is not restrictive and the regression yields the OLS solution; when $s$ is small, the coefficients must be small to avoid budget violation.

\subsection{Geometric Interpretation}

For $p = 2$, the constraint region of the lasso is a \textbf{diamond} ($|\beta_1| + |\beta_2| \leq s$), while the constraint region of ridge regression is a \textbf{disk} ($\beta_1^2 + \beta_2^2 \leq s$). The elliptical contours of the RSS expand outward from the OLS solution $\hat{\beta}$. The regularized solution is the point where the contours first \textbf{touch} the constraint region.

Because the lasso constraint has \textbf{corners} (at the axes), the contours are more likely to intersect the constraint at a corner, where one coefficient is exactly zero. The ridge constraint, being smooth, has no corners, so the intersection almost never occurs at an axis --- hence ridge rarely produces exact zeros.

%-------------------------------------------------------------
\section{Lasso vs Ridge Regression}
%-------------------------------------------------------------

\begin{tcolorbox}[colback=green!5!white,colframe=green!75!black,title=\textbf{Comparison}]
\begin{itemize}
    \item \textbf{Lasso advantage}: produces simpler and more \textbf{interpretable} models that involve only a \textbf{subset of the predictors}, generating more accurate predictions when few predictors are truly related to the response.
    \item \textbf{Ridge advantage}: performs better when the response is a function of \textbf{many predictors}, all with coefficients of roughly equal size.
    \item The number of predictors related to the response is \textbf{never known a priori} in real problems.
    \item \textbf{Cross-validation} is used to determine the better approach on a particular dataset.
\end{itemize}
\end{tcolorbox}

%-------------------------------------------------------------
\section{Selecting the Tuning Parameter}
%-------------------------------------------------------------

Both ridge regression and the lasso require a method for selecting the tuning parameter $\lambda$ (or equivalently the budget $s$).

\begin{tcolorbox}[colback=blue!5!white,colframe=blue!75!black,title=\textbf{Tuning Parameter Selection via Cross-Validation}]
\begin{enumerate}
    \item Choose a grid of $\lambda$ values.
    \item Compute the cross-validation error $\mathrm{CV}_\lambda$ for each $\lambda$.
    \item Select the value of $\lambda$ that minimizes the cross-validation error.
    \item Re-fit the model using all available training data and the selected $\lambda$.
\end{enumerate}
\end{tcolorbox}

The cross-validation error as a function of $\lambda$ is typically a U-shaped curve: too little regularization (small $\lambda$) leads to overfitting (high variance), while too much regularization (large $\lambda$) leads to underfitting (high bias). The optimal $\lambda$ balances the two.


% ========================================
\part{Dimension Reduction}
% ========================================

Subset selection and shrinkage control variance by working with a subset of predictors or by shrinking their coefficients. A complementary family of approaches instead \textbf{transforms the predictors} and then fits a model in the transformed space — this is the idea behind \textbf{dimension reduction}.

% ========================================
\section{Dimension Reduction Methods}
% ========================================

\subsection{The General Framework}

Let $Z_1, Z_2, \ldots, Z_M$ denote $M$ \textbf{linear combinations} of the original $p$ predictors, where $M < p$:
\[
    Z_m = \sum_{j=1}^{p} \phi_{jm} X_j
\]
for some constants $\phi_{1m}, \phi_{2m}, \ldots, \phi_{pm}$ and $m = 1, \ldots, M$. We then fit a linear regression model using least squares on these $M$ transformed predictors:
\[
    y_i = \theta_0 + \sum_{m=1}^{M} \theta_m z_{im}, \qquad i = 1, \ldots, n.
\]

\subsection{Why Dimension Reduction?}

The dimension of the regression problem is \textbf{reduced}: instead of estimating $p + 1$ coefficients $\beta_0, \beta_1, \ldots, \beta_p$, we only need to estimate $M + 1$ coefficients $\theta_0, \theta_1, \ldots, \theta_M$, where $M < p$.

Dimension reduction works in \textbf{two steps}:
\begin{enumerate}
    \item The transformed predictors $Z_1, Z_2, \ldots, Z_M$ are obtained.
    \item A linear model is fit using these $M$ predictors via least squares.
\end{enumerate}

\subsection{Connection to Linear Regression}

The reduced model is a \textbf{special case} of the original linear regression model. Substituting the definition of $Z_m$, one can show:
\[
    \beta_j = \sum_{m=1}^{M} \theta_m \phi_{jm}
\]
This imposes a \textbf{constraint} on the estimated $\beta_j$ coefficients, which introduces a potential bias. However:
\begin{itemize}
    \item When $p$ is large relative to $n$, selecting $M \ll p$ can significantly \textbf{reduce the variance} of the fitted coefficients — a worthwhile bias-variance trade-off.
    \item When $M = p$ and all $Z_m$ are linearly independent, no constraints are imposed and fitting the reduced system is \textbf{equivalent} to fitting the original model with coefficients $\beta_j$.
\end{itemize}

The choice of $Z_1, \ldots, Z_M$ (i.e., the selection of $\phi_{jm}$) can be achieved in different ways. The two approaches covered here are \textbf{principal components} and \textbf{partial least squares}.

% ========================================
\section{Principal Components Analysis}
% ========================================

\subsection{Overview}

\begin{tcolorbox}[colback=blue!5!white,colframe=blue!75!black,title=\textbf{Principal Components Analysis (PCA)}]
PCA is the process by which \textbf{principal components} are computed and used to understand the data. It derives a \textbf{low-dimensional set of features} from a large one, finding a small number of dimensions that are as \textbf{interesting} as possible — where ``interesting'' is measured by the amount that the observations \textbf{vary} along each dimension.
\end{tcolorbox}

Key properties of PCA:
\begin{itemize}
    \item PCA is an \textbf{unsupervised approach}: it involves only the features $X_1, X_2, \ldots, X_p$ and no associated response $Y$.
    \item It is primarily a tool for \textbf{data visualization}, but is also useful in \textbf{supervised learning} (via Principal Components Regression).
    \item Each dimension found by PCA is a \textbf{linear combination} of the $p$ features.
\end{itemize}

\subsection{Principal Component Directions}

Consider $n$ observations with measurements on $p$ features. Each observation lives in a $p$-dimensional space, but \textbf{not all dimensions are equally interesting}. PCA seeks directions along which the data are most variable.

\begin{tcolorbox}[colback=blue!5!white,colframe=blue!75!black,title=\textbf{Principal Component Directions}]
Principal component directions are the directions in feature space along which the original data are \textbf{most variable}.
\end{tcolorbox}

\textit{Example}: population size and ad spending for 100 cities. The \textbf{first principal component direction} (green solid line) captures the direction of maximum variance. The \textbf{second principal component direction} (blue dashed line) is perpendicular to the first and captures the remaining variance.

\subsection{The First Principal Component}

The \textbf{first principal component} of a set of $p$ features is the normalized linear combination of the features:
\[
    Z_1 = \phi_{11} X_1 + \phi_{21} X_2 + \cdots + \phi_{p1} X_p
\]
with the \textbf{largest variance}, subject to the normalization constraint $\sum_{j=1}^{p} \phi_{j1}^2 = 1$.

The elements $\phi_{11}, \ldots, \phi_{p1}$ are the \textbf{loadings} of the first principal component; together they form the \textbf{loading vector} $\phi_1 = (\phi_{11}, \phi_{21}, \ldots, \phi_{p1})^T$.

\subsubsection{Optimization Problem}

We solve:
\[
    \max_{\phi_{11},\ldots,\phi_{p1}} \left\{ \frac{1}{n} \sum_{i=1}^{n} \left( \sum_{j=1}^{p} \phi_{j1} x_{ij} \right)^2 \right\}
    \qquad \text{subject to} \qquad
    \sum_{j=1}^{p} \phi_{j1}^2 = 1
\]

The values $z_{11}, \ldots, z_{n1}$ are known as the \textbf{principal component scores}: $z_{i1} = \sum_{j=1}^p \phi_{j1} x_{ij}$.

\subsubsection{Geometric Interpretation}

The loading vector $\phi_1$ defines the line in $p$-dimensional feature space that is \textbf{as close as possible to the data}, in the sense that it \textbf{minimizes the sum of squared perpendicular distances} between each point and the line. Projecting observations onto this line yields principal component scores.

\textit{Example (advertising data)}: the first principal component is:
\[
    Z_1 = 0.839 \times (\mathit{pop} - \overline{\mathit{pop}}) + 0.544 \times (\mathit{ad} - \overline{\mathit{ad}})
\]
where $\phi_{11} = 0.839$ and $\phi_{21} = 0.544$ are the loadings; $\overline{\mathit{pop}}$ and $\overline{\mathit{ad}}$ are sample means (features are centered before computing PCA).

\subsection{The Second Principal Component}

We can construct up to $p$ distinct principal components. The \textbf{second principal component} $Z_2$ is the linear combination of the variables that:
\begin{itemize}
    \item has \textbf{largest variance} subject to being \textbf{uncorrelated with $Z_1$},
    \item must be \textbf{perpendicular} (orthogonal) to the first principal component direction.
\end{itemize}

For the advertising data (with $p = 2$):
\[
    Z_2 = 0.544 \times (\mathit{pop} - \overline{\mathit{pop}}) + 0.839 \times (\mathit{ad} - \overline{\mathit{ad}})
\]

Since $p = 2$, the two principal components together contain \textbf{all of the information} in \textit{pop} and \textit{ad}.

\begin{tcolorbox}[colback=orange!5!white,colframe=orange!75!black,title=\textbf{Key Observation}]
The second principal component scores $z_{i2}$ are much closer to zero than $z_{i1}$, indicating that the second component captures \textbf{less information}. Plots of $z_{i2}$ versus \textit{pop} and \textit{ad} show weak relationships: only the \textbf{first principal component is needed} to accurately represent both predictors.
\end{tcolorbox}

Additional components, if more predictors were available, would \textbf{successively maximize variance} subject to the constraint of being uncorrelated with all preceding components.

% ========================================
\section{Principal Components Regression}
% ========================================

\subsection{The PCR Approach}

\begin{tcolorbox}[colback=blue!5!white,colframe=blue!75!black,title=\textbf{Principal Components Regression (PCR)}]
PCR constructs the first $M$ principal components $Z_1, \ldots, Z_M$, and then uses them as predictors in a \textbf{linear regression model fit by least squares}.
\end{tcolorbox}

\textbf{Key idea}: often a small number of principal components suffice to explain most of the variability in the data, as well as the relationship with the response.

\textbf{Core assumption}: the directions that show the most variation in $X$ are also the most \textbf{associated with $Y$}. This assumption is not guaranteed to be true, but it is often reasonable and provides good empirical results.

\subsection{Advantages of PCR}

If the assumption above holds:
\begin{itemize}
    \item Fitting least squares to $Z_1, \ldots, Z_M$ yields \textbf{better results} than fitting to $X_1, \ldots, X_p$, since most or all of the information relevant to the response is contained in $Z_1, \ldots, Z_M$.
    \item By estimating only $M \ll p$ coefficients we can \textbf{mitigate overfitting}.
\end{itemize}

\textit{Example}: for the advertising data, the first principal component explains most of the variance in both \textit{pop} and \textit{ad}. A PCR model using this single variable to predict \textit{sales} performs well.

\subsection{Choosing $M$ via Cross-Validation}

In PCR, the number of principal components $M$ is a \textbf{tuning parameter} typically chosen by \textbf{cross-validation}. The cross-validation MSE as a function of $M$ is examined; the optimal $M$ minimizes this error.

\subsection{PCR is Not Feature Selection}

\begin{tcolorbox}[colback=orange!5!white,colframe=orange!75!black,title=\textbf{Important: PCR $\neq$ Feature Selection}]
Even though PCR provides a simple way to perform regression using $M < p$ predictors, it is \textbf{not a feature selection method}. Each of the $M$ principal components used in the regression is a \textbf{linear combination of all $p$ original features}.
\end{tcolorbox}

For instance, the first principal component $Z_1$ was a linear combination of both \textit{pop} and \textit{ad}. In this sense, PCR is more closely related to \textbf{ridge regression} than to the lasso. Ridge regression can in fact be considered as a \textbf{continuous version of PCR}.

\subsection{PCR vs Ridge Regression}

Both PCR and ridge regression perform well on data where the first few principal components capture most of the variation relevant to the response. Empirically:
\begin{itemize}
    \item PCR shows a clear minimum in test MSE at some $M^* < p$, with bias dropping rapidly as $M$ increases.
    \item Ridge regression provides a smooth, continuous regularization; PCR provides a discrete one (by choosing how many components to include).
    \item On data designed to depend on the first five principal components, both PCR (at optimal $M$) and ridge outperform the lasso.
\end{itemize}

% ========================================
\section{Partial Least Squares}
% ========================================

\subsection{Motivation: The Supervised Alternative to PCA}

PCR identifies directions that best represent the predictors $X_1, \ldots, X_p$ in an \textbf{unsupervised} way: the response $Y$ plays no role in determining the principal component directions.

\begin{tcolorbox}[colback=orange!5!white,colframe=orange!75!black,title=\textbf{Main Drawback of PCR}]
There is no guarantee that the directions that best explain the predictors will also be the best directions for \textbf{predicting the response}.
\end{tcolorbox}

\begin{tcolorbox}[colback=blue!5!white,colframe=blue!75!black,title=\textbf{Partial Least Squares (PLS)}]
PLS is a \textbf{supervised} dimension reduction method. It identifies a new set of features $Z_1, \ldots, Z_M$ that are linear combinations of the original features, and then fits a linear model via least squares using these $M$ new features. Unlike PCR, PLS identifies the new features in a \textbf{supervised} way — making use of $Y$ to find directions that explain both the predictors and the response.
\end{tcolorbox}

\subsection{Computing PLS Directions}

\textbf{First PLS direction $Z_1$}: after standardizing the $p$ predictors, PLS computes the first direction by setting each loading $\phi_{j1}$ equal to the \textbf{coefficient from the simple linear regression of $Y$ onto $X_j$}. This coefficient is proportional to the correlation between $Y$ and $X_j$. Hence:
\[
    Z_1 = \sum_{j=1}^{p} \phi_{j1} X_j
\]
places the \textbf{highest weight} on the variables most strongly correlated with the response.

\textit{Example}: on the advertising data, the first PLS direction has less change in the \textit{ad} dimension per unit change in the \textit{pop} dimension, relative to PCA. This suggests that \textit{pop} is more highly correlated with the response than \textit{ad}.

\textbf{Second PLS direction $Z_2$}: we first \textbf{adjust} each variable for $Z_1$, by regressing each variable on $Z_1$ and taking the \textbf{residuals} (the information not explained by the first PLS direction). We then compute $Z_2$ from this orthogonalized data in the same fashion as $Z_1$.

This iterative procedure can be repeated up to $M$ times to identify PLS components $Z_1, \ldots, Z_M$.

\subsection{Fitting and Tuning}

At the end of the procedure, least squares is used to fit a linear model to predict $Y$ using $Z_1, \ldots, Z_M$. As with PCR, the \textbf{number $M$ of PLS directions is a tuning parameter} typically chosen by \textbf{cross-validation}.

\begin{remark}
While PLS can reduce bias relative to PCR (since it uses $Y$ to find directions), it can also \textbf{increase variance}. In practice, PLS often performs no better than ridge regression or PCR. The theoretical properties of PLS are less well understood than those of PCR.
\end{remark}

% ========================================
\section{Considerations in High Dimensions}
% ========================================

\subsection{Low-Dimensional vs High-Dimensional Data}

Most traditional statistical techniques for regression and classification are intended for the \textbf{low-dimensional setting} in which $n \gg p$ (many more observations than features).

\textit{Example}: predicting a patient's blood pressure from age, gender, and BMI. Here $p = 3$ (or 4 with an intercept) and $n$ could be several thousand patients — clearly $n \gg p$, a low-dimensional problem.

In the past 20 years, new technologies have enabled collection of an almost unlimited number of feature measurements ($p$ very large), while $n$ may be limited by cost or sample availability. Data sets with \textbf{more features than observations} ($p > n$) are called \textbf{high-dimensional}.

\textit{Examples}:
\begin{itemize}
    \item Predicting blood pressure from half a million SNPs (single nucleotide polymorphisms): $n \approx 200$, $p \approx 500{,}000$.
    \item Online shopping pattern analysis using the \textbf{bag-of-words} model: each search term is a binary feature; $p$ can easily exceed $n$.
\end{itemize}

\subsection{What Goes Wrong in High Dimensions?}

Classical approaches such as \textbf{least squares linear regression are not appropriate} in the high-dimensional setting. When $p \geq n$, least squares will yield coefficient estimates that result in a \textbf{perfect fit to the training data} (residuals are exactly zero) — regardless of whether there truly is a relationship between the features and the response.

\begin{tcolorbox}[colback=orange!5!white,colframe=orange!75!black,title=\textbf{The High-Dimensional Danger}]
When $p \geq n$, least squares always produces a perfect fit to training data. This perfect fit \textbf{almost certainly leads to overfitting}: the model performs extremely poorly on an independent test set, since it is too flexible and captures noise.
\end{tcolorbox}

\textit{Illustration}: with $n = 20$ observations and $p = 1$ predictor, least squares does not perfectly fit the data. With only $n = 2$ observations, any line passes through both points — a perfect fit that generalizes nothing.

\subsection{The Curse of Dimensionality in Practice}

The principal issues in high-dimensional analysis are:
\begin{itemize}
    \item The \textbf{bias-variance trade-off} becomes critical: adding more features can increase variance dramatically.
    \item The danger of \textbf{overfitting} is amplified when $p$ is large relative to $n$.
    \item These considerations also apply when $p$ is only \textbf{slightly smaller} than $n$.
\end{itemize}

Empirically: simulating $n = 20$ observations with regression on $p$ varying from 1 to 20 (all features unrelated to the response):
\begin{itemize}
    \item $R^2$ \textbf{increases to 1} as $p$ increases — even though no feature is relevant.
    \item Training MSE \textbf{decreases to 0} as $p$ increases.
    \item Test MSE \textbf{increases dramatically} as $p$ increases.
\end{itemize}

Including additional predictors leads to a \textbf{huge increase in the variance of the coefficient estimates}. This highlights the importance of always evaluating model performance on an \textbf{independent test set}, and of \textbf{never relying solely on $R^2$ or training MSE} as measures of model quality.

\subsection{Implications for Model Selection Criteria}

Approaches for adjusting training RSS or $R^2$ to account for model size — such as $C_p$, AIC, and BIC — are \textbf{not appropriate in the high-dimensional setting}, because estimating $\hat{\sigma}^2$ is problematic when $p \geq n$ (the residual degrees of freedom become zero or negative).

\begin{tcolorbox}[colback=green!5!white,colframe=green!75!black,title=\textbf{High-Dimensional Takeaways}]
\begin{itemize}
    \item Regularization and dimension reduction methods (ridge, lasso, PCR, PLS) are essential in high dimensions.
    \item Always report performance on an \textbf{independent test set} or via \textbf{cross-validation}.
    \item Be cautious of $R^2$ and training MSE — they are unreliable indicators of model quality when $p$ is large.
    \item The best model in high dimensions contains \textbf{only a small number of relevant features}: use regularization to enforce sparsity.
\end{itemize}
\end{tcolorbox}

% ========================================
\part{Classification}
% ========================================

Classification is the task of predicting a \textbf{qualitative} (categorical) response. Just as regression is used for quantitative outputs, classification methods are used when the output takes on one of $K$ discrete classes. Many of the core ideas from regression — bias-variance trade-off, training vs.\ test error, model selection — carry over with only minor modifications.

% ========================================
\section{Introduction to Classification}
% ========================================

\subsection{Qualitative Variables}

In many situations, the response variable is \textbf{qualitative}: it takes values from a finite, unordered set of categories (e.g., eye color: blue, brown, green). The process of predicting a qualitative response is known as \textbf{classification} — it involves \textbf{assigning the observation to a category, or class}.

Classifiers typically work by first \textbf{estimating the probability} that an observation belongs to each possible class, and then making a decision based on those probabilities. Three of the most widely used classifiers are:
\begin{itemize}
    \item \textbf{K-Nearest Neighbours (KNN)}
    \item \textbf{Logistic Regression}
    \item \textbf{Linear Discriminant Analysis (LDA)}
\end{itemize}

\textit{Example applications}:
\begin{enumerate}
    \item A patient arrives at the emergency room with a set of symptoms — classify into one of three medical conditions.
    \item Fraud detection: determine whether a transaction is fraudulent, based on IP address, transaction history, etc.
    \item Predict, on the basis of DNA sequence data, which mutations are disease-causing.
\end{enumerate}

\subsection{The Classification Setting}

We seek to estimate the classification function $f$ from training observations $\{(x_1, y_1), \ldots, (x_n, y_n)\}$ where $y_1, \ldots, y_n$ are qualitative. The accuracy of $\hat{f}$ is quantified by the \textbf{training error rate}:
\[
    \varepsilon_{\text{train}} = \frac{1}{n} \sum_{i=1}^{n} I(y_i \neq \hat{y}_i)
\]
where $\hat{y}_i$ is the predicted class label and $I(\cdot)$ is the indicator function (1 if $y_i \neq \hat{y}_i$, zero otherwise). $\varepsilon_{\text{train}}$ is simply the fraction of incorrect classifications on the training data. The \textbf{test error rate} is computed from predictions on data not used for training — this is what we actually care about.

\subsection{The Bayes Classifier}

\begin{tcolorbox}[colback=blue!5!white,colframe=blue!75!black,title=\textbf{The Bayes Classifier}]
The test error rate is minimized, on average, by a classifier that assigns each observation to the \textbf{most likely class given its predictor values}. That is, assign a test observation with predictor vector $x_0$ to the class $j$ for which
\[
    \Pr(Y = j \mid X = x_0)
\]
is the \textbf{largest}. This is the \textbf{Bayes classifier}.
\end{tcolorbox}

In a \textbf{two-class problem} (classes 1 and 2), the Bayes classifier assigns an observation to class 1 if $\Pr(Y = 1 \mid X = x_0) > 0.5$, and to class 2 otherwise.

\subsubsection{Bayes Decision Boundary}

The \textbf{Bayes decision boundary} is the set of points where the largest conditional class probability is exactly 0.5 (in the two-class case). Observations on one side of the boundary are assigned to one class, and observations on the other side to the other class.

The Bayes classifier produces the \textbf{lowest possible test error rate}, called the \textbf{Bayes error rate}. At a given $X = x_0$, the error rate is $1 - \max_j \Pr(Y = j \mid X = x_0)$. The overall Bayes error rate is:
\[
    1 - \mathbb{E}\!\left[\max_j \Pr(Y = j \mid X)\right]
\]
where the expectation averages over all possible values of $X$.

\begin{tcolorbox}[colback=orange!5!white,colframe=orange!75!black,title=\textbf{The Bayes Classifier is Unattainable}]
For real data, we do \textbf{not} know the conditional distribution of $Y$ given $X$, so computing the Bayes classifier is \textbf{impossible}. It serves as an \textbf{unattainable gold standard} against which to compare other methods.
\end{tcolorbox}

% ========================================
\section{K-Nearest Neighbours}
% ========================================

\subsection{The KNN Classifier}

Many methods estimate the conditional distribution of $Y$ given $X$, and then classify an observation to the class with the highest estimated probability. One of the most widely used is the \textbf{K-Nearest Neighbours (KNN)} classifier.

\begin{tcolorbox}[colback=blue!5!white,colframe=blue!75!black,title=\textbf{KNN Algorithm}]
Given a positive integer $K$ and a test observation $x_0$, KNN:
\begin{enumerate}
    \item Identifies the $K$ training points closest to $x_0$ — call this set $\mathcal{N}_0$.
    \item Estimates the conditional probability for class $j$ as the \textbf{fraction of points in $\mathcal{N}_0$ whose response equals $j$}:
    \[
        \Pr(Y = j \mid X = x_0) = \frac{1}{K} \sum_{i \in \mathcal{N}_0} I(y_i = j)
    \]
    \item Classifies $x_0$ to the class with the \textbf{largest probability}.
\end{enumerate}
\end{tcolorbox}

\textit{Example with $K = 3$}: to predict the class of a test point (black cross), KNN finds the three closest training observations — suppose two are blue and one is orange. The estimated probabilities are $2/3$ for blue and $1/3$ for orange. KNN predicts the \textbf{blue class}.

\subsection{Choice of $K$: Bias-Variance Trade-off}

The value of $K$ is the key tuning parameter and controls the \textbf{flexibility} of the classifier:

\begin{itemize}
    \item \textbf{$K = 1$}: the decision boundary is \textbf{overly flexible}, fits patterns not corresponding to the Bayes boundary — \textbf{low bias, very high variance}. Training error rate is exactly 0.
    \item \textbf{Large $K$} (e.g., $K = 100$): the boundary becomes \textbf{less flexible}, close to linear — \textbf{lower variance, higher bias}.
    \item \textbf{Optimal $K$} (e.g., $K = 10$ on a simulated dataset): KNN decision boundary is close to the Bayes boundary; test error rate (0.1363) is close to the Bayes error rate (0.1304).
\end{itemize}

\begin{tcolorbox}[colback=blue!5!white,colframe=blue!75!black,title=\textbf{Flexibility and Error in KNN}]
As $1/K$ increases (more flexible), training error consistently decreases. However, the test error forms a characteristic \textbf{U-shape}: too flexible (small $K$) leads to overfitting and high test error; too rigid (large $K$) leads to underfitting. The optimal $K$ minimizes test error.
\end{tcolorbox}

% ========================================
\part{Classification II --- Logistic Regression}
% ========================================

% ========================================
\section{Why Not Linear Regression?}
% ========================================

\subsection{The Problem with Encoding Qualitative Responses}

A natural first attempt at classification is to encode the classes as numbers and apply linear regression. However, this approach has a fundamental flaw: the \textbf{order of the encoding matters}.

\textit{Example}: predict a medical condition — stroke, drug overdose, epileptic seizure. We could encode:
\[
    Y = \begin{cases} 1 & \text{if stroke} \\ 2 & \text{if drug overdose} \\ 3 & \text{if epileptic seizure} \end{cases}
\]
and use least squares to fit a linear model on the predictors $X_1, \ldots, X_p$. But this encoding \textbf{implies an ordering and equal spacing} between conditions that does not exist medically. A different encoding (e.g., stroke = 1, epileptic seizure = 2, drug overdose = 3) would produce \textbf{fundamentally different linear models} and hence different predictions.

The only situation where a quantitative encoding is defensible is when the response classes have a \textbf{natural ordering and equal gaps} — for instance:
\[
    Y = \begin{cases} 1 & \text{if mild} \\ 2 & \text{if moderate} \\ 3 & \text{if severe} \end{cases}
\]
In general, there is \textbf{no natural way} to encode a qualitative response with more than two values into a single quantitative response.

\subsection{The Binary Case: Linear Regression on a Dummy Variable}

When the response has only \textbf{two classes}, we can use a binary (dummy variable) encoding:
\[
    Y = \begin{cases} 0 & \text{if stroke} \\ 1 & \text{if drug overdose} \end{cases}
\]
and fit a linear regression model to predict $\Pr(\text{drug overdose} \mid X) \approx X\hat{\beta}$. Then classify as drug overdose if $\hat{Y} > 0.5$, and stroke otherwise.

This works for two classes, but has two critical issues:
\begin{enumerate}
    \item The linear regression estimate $X\hat{\beta}$ can fall \textbf{outside the interval $[0,1]$} — giving negative probabilities or probabilities above 1, which are not valid.
    \item The approach \textbf{cannot be extended} to qualitative responses with more than two values.
\end{enumerate}

\begin{tcolorbox}[colback=orange!5!white,colframe=orange!75!black,title=\textbf{Conclusion}]
We need classification methods that are specifically designed for qualitative responses and that always produce \textbf{valid probability estimates in $[0, 1]$}.
\end{tcolorbox}

% ========================================
\section{Logistic Regression}
% ========================================

\subsection{The Logistic Model}

Rather than modelling $Y$ directly, logistic regression models the \textbf{probability that $Y$ belongs to a particular category}. For the Default dataset — response: default (Yes/No), predictor: credit card balance — we model:
\[
    p(\text{balance}) = \Pr(\text{default} = \text{Yes} \mid \text{balance})
\]

Using a linear model $p(X) = \beta_0 + \beta_1 X$ leads to estimates outside $[0,1]$ for extreme values of $X$. To avoid this, logistic regression uses the \textbf{logistic function}:

\begin{tcolorbox}[colback=blue!5!white,colframe=blue!75!black,title=\textbf{The Logistic Function}]
\[
    p(X) = \frac{e^{\beta_0 + \beta_1 X}}{1 + e^{\beta_0 + \beta_1 X}}
\]
\end{tcolorbox}

The logistic function always produces outputs in $(0, 1)$ for any value of $X$, giving an \textbf{S-shaped curve} — always a reasonable probability prediction. The model is fit using \textbf{maximum likelihood}.

\textit{Illustration}: for the Default data, linear regression produces some \textbf{negative} estimated probabilities of default. Logistic regression keeps all probabilities in $[0,1]$.

\subsection{The Odds}

After algebraic manipulation of the logistic function, we obtain:
\[
    \frac{p(X)}{1 - p(X)} = e^{\beta_0 + \beta_1 X}
\]

The quantity $p(X) / (1 - p(X))$ is called the \textbf{odds}, and ranges in $[0, \infty)$:
\begin{itemize}
    \item Odds close to 0: very low probability of the event.
    \item Odds much larger than 1: very high probability.
    \item \textit{Example}: $p(X) = 0.2 \Rightarrow$ odds $= 0.2/0.8 = 1/4$ (on average 1 in 5 people will default).
    \item \textit{Example}: $p(X) = 0.9 \Rightarrow$ odds $= 0.9/0.1 = 9$ (9 out of 10 will default).
\end{itemize}

\subsection{The Logit (Log-Odds)}

Taking the \textbf{logarithm} of both sides:
\[
    \log\!\left(\frac{p(X)}{1 - p(X)}\right) = \beta_0 + \beta_1 X
\]

The left-hand side is called the \textbf{logit} or \textbf{log-odds}. Key properties:
\begin{itemize}
    \item The logit is \textbf{linear in $X$}: the logistic regression model is linear on the log-odds scale.
    \item Increasing $X$ by one unit changes the log odds by $\beta_1$, equivalently multiplying the odds by $e^{\beta_1}$.
    \item The change in $p(X)$ per unit increase in $X$ \textbf{depends on the current value of $X$} (non-constant, unlike linear regression).
    \item If $\beta_1 > 0$: increasing $X$ \textbf{increases} $p(X)$. If $\beta_1 < 0$: increasing $X$ \textbf{decreases} $p(X)$.
\end{itemize}

\subsection{Estimating the Coefficients}

The coefficients $\beta_0$ and $\beta_1$ are estimated from training data. We could use (non-linear) least squares, but \textbf{maximum likelihood} is preferred as it has better statistical properties.

\begin{tcolorbox}[colback=blue!5!white,colframe=blue!75!black,title=\textbf{The Likelihood Function}]
We seek $\hat{\beta}_0$ and $\hat{\beta}_1$ such that predicted probabilities are \textbf{close to 1} for individuals who defaulted and \textbf{close to 0} for those who did not. This is formalized by the likelihood function:
\[
    \ell(\beta_0, \beta_1) = \prod_{i:\, y_i = 1} p(x_i) \prod_{i':\, y_{i'} = 0} \bigl(1 - p(x_{i'})\bigr)
\]
The estimates $\hat{\beta}_0$ and $\hat{\beta}_1$ are chosen to \textbf{maximize} $\ell(\beta_0, \beta_1)$.
\end{tcolorbox}

\begin{remark}
Least squares in the linear regression setting is a special case of maximum likelihood. Logistic regression and other models can typically be fit using a standard statistical software package.
\end{remark}

\subsubsection{Coefficient Estimates for the Default Dataset}

Fitting logistic regression on the Default data to predict $\Pr(\text{default} = \text{Yes} \mid \text{balance})$:

\begin{center}
\begin{tabular}{lcccc}
\hline
 & Coefficient & Std.\ error & Z-statistic & P-value \\
\hline
Intercept & $-10.6513$ & $0.3612$ & $-29.5$ & $<0.0001$ \\
balance   & $0.0055$   & $0.0002$ & $24.9$  & $<0.0001$ \\
\hline
\end{tabular}
\end{center}

$\hat{\beta}_1 = 0.0055$: a one-unit increase in balance increases the \textbf{log odds of default by 0.0055}.

Many aspects of logistic regression are \textbf{analogous to linear regression}:
\begin{itemize}
    \item The \textbf{z-statistic} plays the role of the t-statistic: $z = \hat{\beta}_1 / \mathrm{SE}(\hat{\beta}_1)$. A large absolute value is evidence against $H_0: \beta_1 = 0$.
    \item The \textbf{p-value} tests whether there is a statistically significant association between balance and default.
\end{itemize}

\subsubsection{Making Predictions}

Once coefficients are estimated, computing probabilities is straightforward. For a balance of \$1000:
\[
    \hat{p}(X) = \frac{e^{-10.6513 + 0.0055 \times 1000}}{1 + e^{-10.6513 + 0.0055 \times 1000}} \approx 0.00576
\]
which is below 1\%. For a balance of \$2000 the predicted probability is $0.586 = 58.6\%$ — much higher.

\subsubsection{Using Qualitative Predictors}

Qualitative predictors are included via \textbf{dummy variables}. For the Default dataset, student status is encoded as 0 (non-student) and 1 (student):

\begin{center}
\begin{tabular}{lcccc}
\hline
 & Coefficient & Std.\ error & z-statistic & p-value \\
\hline
Intercept      & $-3.5041$ & $0.0707$ & $-49.55$ & $<0.0001$ \\
student[Yes]   & $0.4049$  & $0.1150$ & $3.52$   & $0.0004$  \\
\hline
\end{tabular}
\end{center}

The \textbf{positive coefficient} for student[Yes] indicates that students have a \textbf{higher chance of default} than non-students (when using student status as the only predictor).

% ========================================
\section{Multiple Logistic Regression}
% ========================================

\subsection{Extending to Multiple Predictors}

By analogy with the extension from simple to multiple linear regression, the logistic model generalizes to \textbf{multiple predictors} $X = (X_1, \ldots, X_p)$:

\begin{tcolorbox}[colback=blue!5!white,colframe=blue!75!black,title=\textbf{Multiple Logistic Regression}]
\[
    \log\!\left(\frac{p(X)}{1 - p(X)}\right) = \beta_0 + \beta_1 X_1 + \cdots + \beta_p X_p
\]
equivalently:
\[
    p(X) = \frac{e^{\beta_0 + \beta_1 X_1 + \cdots + \beta_p X_p}}{1 + e^{\beta_0 + \beta_1 X_1 + \cdots + \beta_p X_p}}
\]
Parameters $\beta_0, \beta_1, \ldots, \beta_p$ are estimated via \textbf{maximum likelihood}.
\end{tcolorbox}

\subsection{Example: Default Dataset with Three Predictors}

Fitting logistic regression using balance, income, and student status simultaneously:

\begin{center}
\begin{tabular}{lcccc}
\hline
 & Coefficient & Std.\ error & z-statistic & p-value \\
\hline
Intercept      & $-10.8690$ & $0.4923$ & $-22.08$ & $<0.0001$ \\
balance        & $0.0057$   & $0.0002$ & $24.74$  & $<0.0001$ \\
income         & $0.0030$   & $0.0082$ & $0.37$   & $0.7115$  \\
student[Yes]   & $-0.6468$  & $0.2362$ & $-2.74$  & $0.0062$  \\
\hline
\end{tabular}
\end{center}

\begin{tcolorbox}[colback=orange!5!white,colframe=orange!75!black,title=\textbf{Confounding: the Student Sign Reversal}]
When using student status alone, the coefficient is \textbf{positive} (+0.4049): students appear to default more. When controlling for balance and income, the coefficient becomes \textbf{negative} ($-0.6468$): students actually default \textit{less} at a fixed balance level.

This is a case of \textbf{confounding}: students tend to carry higher credit card balances, and high balance is the true driver of default. Once balance is held fixed, being a student is associated with \textbf{lower} default risk. This illustrates the importance of considering all relevant predictors simultaneously — a single-predictor model can give a misleading picture.
\end{tcolorbox}

% ========================================
\part{Classification III --- Linear Discriminant Analysis}
% ========================================

% ========================================
\section{Introduction}
% ========================================

\subsection{An Alternative Approach to Classification}

Logistic regression models $\Pr(Y = k \mid X = x)$ \textbf{directly} using the logistic function. Linear Discriminant Analysis (LDA) takes a different, \textbf{less direct} approach:
\begin{enumerate}
    \item Model the distribution of the predictors $X$ \textbf{separately within each class} — i.e., estimate $f_k(x) = \Pr(X = x \mid Y = k)$ for each class $k$.
    \item Use \textbf{Bayes' theorem} to flip these around into estimates of $\Pr(Y = k \mid X = x)$.
\end{enumerate}

When the within-class distributions are assumed to be \textbf{Gaussian}, the resulting model turns out to be very similar to logistic regression.

\subsection{Why Use LDA?}

\begin{tcolorbox}[colback=blue!5!white,colframe=blue!75!black,title=\textbf{Advantages of LDA over Logistic Regression}]
\begin{itemize}
    \item When the classes are \textbf{well-separated}, parameter estimates for logistic regression are surprisingly \textbf{unstable}. LDA does not suffer from this problem.
    \item When $n$ is \textbf{small} and the distribution of $X$ is approximately \textbf{normal} within each class, LDA is more stable than logistic regression.
    \item LDA is naturally suited for problems with \textbf{more than two response classes}.
\end{itemize}
\end{tcolorbox}

% ========================================
\section{Bayes' Theorem for Classification}
% ========================================

\subsection{Setup}

Suppose we wish to classify an observation into one of $K$ classes, where $K \geq 2$ — the categorical variable $Y$ takes $K$ possible distinct and unordered values. Define:
\begin{itemize}
    \item $\pi_k$ = \textbf{prior probability} that a randomly chosen observation comes from class $k$; this is the probability that a given observation belongs to the $k$-th category of $Y$.
    \item $f_k(x) = \Pr(X = x \mid Y = k)$ = \textbf{density function} of $X$ for an observation from class $k$. $f_k(x)$ is large if it is highly probable that an observation in class $k$ has $X \approx x$, and small if it is very unlikely.
\end{itemize}

\subsection{Bayes' Theorem}

\begin{tcolorbox}[colback=blue!5!white,colframe=blue!75!black,title=\textbf{Bayes' Theorem for Classification}]
\[
    p_k(x) = \Pr(Y = k \mid X = x) = \frac{\pi_k f_k(x)}{\displaystyle\sum_{j=1}^{K} \pi_j f_j(x)}
\]
\end{tcolorbox}

$p_k(x)$ is the \textbf{posterior probability} that an observation $X = x$ belongs to class $k$ — i.e., the probability that the observation belongs to the $k$-th class, given the predictor value.

\begin{itemize}
    \item Estimating $\pi_k$ is \textbf{easy}: compute the \textbf{fraction of training observations} that belong to class $k$.
    \item Estimating $f_k(x)$ is more \textbf{challenging} — we need to assume some form for these densities. LDA assumes they are Gaussian.
\end{itemize}

The Bayes classifier, which achieves the lowest possible error rate, assigns $X = x$ to the class with the largest $p_k(x)$. If we can estimate $f_k(x)$, we can build a classifier that \textbf{approximates the Bayes classifier}.

% ========================================
\section{Linear Discriminant Analysis for \texorpdfstring{$p = 1$}{p=1}}
% ========================================

\subsection{Gaussian Assumption}

With a single predictor ($p = 1$), assume $f_k(x)$ is \textbf{normal (Gaussian)}:
\[
    f_k(x) = \frac{1}{\sqrt{2\pi}\,\sigma_k} \exp\!\left(-\frac{1}{2\sigma_k^2}(x - \mu_k)^2\right)
\]
where $\mu_k$ and $\sigma_k^2$ are the mean and variance for class $k$.

LDA further assumes a \textbf{shared variance} across all $K$ classes: $\sigma_1^2 = \cdots = \sigma_K^2 = \sigma^2$.

Plugging into the Bayes formula:
\[
    p_k(x) = \frac{\pi_k \dfrac{1}{\sqrt{2\pi}\,\sigma} \exp\!\left(-\dfrac{1}{2\sigma^2}(x - \mu_k)^2\right)}{\displaystyle\sum_{j=1}^{K} \pi_j \dfrac{1}{\sqrt{2\pi}\,\sigma} \exp\!\left(-\dfrac{1}{2\sigma^2}(x - \mu_j)^2\right)}
\]

The Bayes classifier assigns $X = x$ to the class for which this expression is largest.

\subsection{Discriminant Functions}

Taking the logarithm and rearranging, maximizing $p_k(x)$ over $k$ is equivalent to assigning $x$ to the class with the largest \textbf{linear discriminant function}:

\begin{tcolorbox}[colback=blue!5!white,colframe=blue!75!black,title=\textbf{Linear Discriminant Function ($p = 1$)}]
\[
    \delta_k(x) = x \cdot \frac{\mu_k}{\sigma^2} - \frac{\mu_k^2}{2\sigma^2} + \log(\pi_k)
\]
\end{tcolorbox}

These functions are \textbf{linear in $x$} — hence the name \textit{Linear} Discriminant Analysis.

\textit{Example ($K = 2$, $\pi_1 = \pi_2$)}: the Bayes classifier assigns $x$ to class 1 if $2x(\mu_1 - \mu_2) > \mu_1^2 - \mu_2^2$, and to class 2 otherwise. The \textbf{Bayes decision boundary} is the point where $\delta_1(x) = \delta_2(x)$:
\[
    x = \frac{\mu_1^2 - \mu_2^2}{2(\mu_1 - \mu_2)} = \frac{\mu_1 + \mu_2}{2}
\]

\subsection{The LDA Algorithm}

In practice, $\mu_k$, $\pi_k$, and $\sigma^2$ are unknown. LDA estimates them from training data:

\begin{tcolorbox}[colback=blue!5!white,colframe=blue!75!black,title=\textbf{LDA Parameter Estimates}]
\[
    \hat{\mu}_k = \frac{1}{n_k} \sum_{i:\, y_i = k} x_i, \qquad
    \hat{\sigma}^2 = \frac{1}{n - K} \sum_{k=1}^{K} \sum_{i:\, y_i = k} (x_i - \hat{\mu}_k)^2, \qquad
    \hat{\pi}_k = \frac{n_k}{n}
\]
where $n$ is the total number of training observations and $n_k$ is the number in class $k$.
\end{tcolorbox}

\begin{itemize}
    \item $\hat{\mu}_k$: average of all training observations from class $k$.
    \item $\hat{\sigma}^2$: weighted average of the sample variances for each class.
    \item $\hat{\pi}_k$: proportion of training observations in class $k$ (or the known prior, if available).
\end{itemize}

LDA plugs these estimates into $\delta_k(x)$ and assigns $X = x$ to the class for which:
\[
    \hat{\delta}_k(x) = x \cdot \frac{\hat{\mu}_k}{\hat{\sigma}^2} - \frac{\hat{\mu}_k^2}{2\hat{\sigma}^2} + \log(\hat{\pi}_k)
\]
is the largest.

\subsection{Decision Boundary and Performance}

\textit{Example}: 20 observations drawn from each of two Gaussian classes ($\mu_1 = -1.25$, $\mu_2 = 1.25$, $\sigma^2 = 1$). The two density functions overlap, so there is inherent uncertainty about class membership.

Since $n_1 = n_2 = 20$, we have $\hat{\pi}_1 = \hat{\pi}_2$. The \textbf{LDA decision boundary} is the midpoint of the sample means:
\[
    \frac{\hat{\mu}_1 + \hat{\mu}_2}{2}
\]
The \textbf{Bayes decision boundary} (optimal) is at $(\mu_1 + \mu_2)/2 = 0$. The LDA boundary (solid line) is close to but not exactly at the Bayes boundary (dashed line), due to estimation from finite data.

\begin{tcolorbox}[colback=green!5!white,colframe=green!75!black,title=\textbf{LDA Performance}]
\begin{itemize}
    \item Bayes error rate: \textbf{10.6\%} (theoretical minimum).
    \item LDA test error rate: \textbf{11.1\%} — only 0.5\% above optimal.
\end{itemize}
LDA closely approximates the Bayes classifier when the Gaussian assumption is satisfied.
\end{tcolorbox}

% ========================================
\section{Linear Discriminant Analysis for \texorpdfstring{$p > 1$}{p>1}}
% ========================================

\subsection{Multivariate Gaussian Distribution}

For multiple predictors, LDA assumes $X = (X_1, \ldots, X_p)$ follows a \textbf{multivariate Gaussian distribution} $X \sim \mathcal{N}(\mu, \Sigma)$:
\begin{itemize}
    \item Each individual predictor follows a one-dimensional normal distribution, with some \textbf{correlation} between each pair of predictors.
    \item $E(X) = \mu$ is a vector with $p$ components (the mean of $X$).
    \item $\mathrm{Cov}(X) = \Sigma$ is the $p \times p$ \textbf{covariance matrix} of $X$.
\end{itemize}

The multivariate Gaussian density is:
\[
    f(x) = \frac{1}{(2\pi)^{p/2} |\Sigma|^{1/2}} \exp\!\left(-\frac{1}{2}(x - \mu)^T \Sigma^{-1}(x - \mu)\right)
\]

LDA assumes each class $k$ has its own \textbf{class-specific mean vector} $\mu_k$, but all classes \textbf{share a common covariance matrix} $\Sigma$.

\subsection{Discriminant Functions for \texorpdfstring{$p > 1$}{p>1}}

Plugging in the multivariate Gaussian density, the Bayes classifier assigns $X = x$ to the class for which the \textbf{multivariate linear discriminant function} is largest:

\begin{tcolorbox}[colback=blue!5!white,colframe=blue!75!black,title=\textbf{Linear Discriminant Function ($p > 1$)}]
\[
    \delta_k(x) = x^T \Sigma^{-1} \mu_k - \frac{1}{2} \mu_k^T \Sigma^{-1} \mu_k + \log \pi_k
\]
\end{tcolorbox}

This is the vector/matrix generalization of the univariate formula. Note that $\delta_k(x)$ is a \textbf{linear function of $x$}: the LDA decision rule depends on $x$ only through a \textbf{linear combination} of its elements.

LDA estimates the unknown parameters $\mu_1, \ldots, \mu_K$, $\pi_1, \ldots, \pi_K$, and $\Sigma$ from training data, plugs them into $\delta_k(x)$, and classifies to the class for which $\hat{\delta}_k(x)$ is the largest.

\subsection{Example: Two Predictors, Three Classes}

Three equally-sized Gaussian classes with class-specific mean vectors and a common covariance matrix. The \textbf{Bayes decision boundaries} (dashed lines) are defined by the set of values $x$ for which $\delta_k(x) = \delta_l(x)$, i.e.:
\[
    x^T \Sigma^{-1} \mu_k - \frac{1}{2}\mu_k^T \Sigma^{-1} \mu_k = x^T \Sigma^{-1} \mu_l - \frac{1}{2}\mu_l^T \Sigma^{-1} \mu_l \qquad \text{for } k \neq l
\]

With $K = 3$ classes there are three pairs, hence \textbf{three decision boundaries} dividing the predictor space into three regions. The Bayes classifier assigns an observation to the region in which it is located.

With 20 observations drawn from each class, the \textbf{LDA decision boundaries} (solid lines) are close to the Bayes decision boundaries (dashed lines):
\begin{itemize}
    \item Bayes test error rate: \textbf{7.46\%}
    \item LDA test error rate: \textbf{7.70\%}
\end{itemize}
LDA performs well on this data.

% ========================================
\section{Evaluating Classification Performance}
% ========================================

\subsection{Beyond the Overall Error Rate}

A low overall error rate does not tell the full story. Consider LDA applied to the \textbf{Default dataset} (10,000 training samples, binary response: default Yes/No):
\begin{itemize}
    \item LDA training error rate: \textbf{2.75\%} — appears excellent.
    \item But the \textbf{null classifier} (always predicts ``No default'') achieves \textbf{3.33\%} — barely worse.
    \item Only 3.33\% of individuals actually defaulted, so a useless classifier is almost as good by the overall error rate metric alone.
\end{itemize}

A binary classifier can make two types of error:
\begin{itemize}
    \item \textbf{False Negative (FN)}: a true defaulter incorrectly assigned to the ``No default'' class.
    \item \textbf{False Positive (FP)}: a non-defaulter incorrectly assigned to the ``Default'' class.
\end{itemize}

\subsection{The Confusion Matrix}

A \textbf{confusion matrix} displays the counts of all four prediction outcomes:

\begin{center}
\begin{tabular}{llccc}
\hline
 & & \multicolumn{3}{c}{True default status} \\
 & & No & Yes & Total \\
\hline
\multirow{2}{*}{Predicted default status} & No  & 9\,644 & 252 & 9\,896 \\
                                           & Yes & 23     & 81  & 104    \\
\hline
 & Total & 9\,667 & 333 & 10\,000 \\
\hline
\end{tabular}
\end{center}

Reading the matrix: LDA predicted 104 people would default. Of these, 81 actually defaulted and 23 did not. Of the 333 true defaulters, \textbf{252 (75.7\%) were missed} by LDA — a very high miss rate.

\subsection{Key Metrics}

\begin{center}
\begin{tabular}{llll}
\hline
Name & Definition & Synonyms \\
\hline
False Positive rate  & FP / N      & Type I error, 1 $-$ Specificity \\
True Positive rate   & TP / P      & 1 $-$ Type II error, Sensitivity, Recall \\
Positive Pred.\ value & TP / $P^*$ & Precision \\
Negative Pred.\ value & TN / $N^*$ & --- \\
\hline
\end{tabular}
\end{center}

For the Default example:
\begin{itemize}
    \item \textbf{Sensitivity} (true defaulters correctly identified): $81/333 = 24.3\%$ — very low.
    \item \textbf{Specificity} (non-defaulters correctly identified): $9{,}644/9{,}667 = 99.8\%$ — very high.
\end{itemize}

\begin{tcolorbox}[colback=orange!5!white,colframe=orange!75!black,title=\textbf{Why Does LDA Miss So Many Defaulters?}]
LDA approximates the Bayes classifier, which minimises the \textbf{total} error rate irrespective of which class the errors come from. Since only 3.33\% of observations default, the Bayes classifier sacrifices sensitivity on the minority class (defaulters) to keep the overall error low. A credit card company, however, may specifically want to \textbf{avoid missing defaulters}.
\end{tcolorbox}

\subsection{Adjusting the Decision Threshold}

The Bayes classifier (and by extension LDA) uses a default threshold of \textbf{50\%}:
\[
    \Pr(\text{default} = \text{yes} \mid X = x) > 0.5 \implies \text{predict Default}
\]

To \textbf{increase sensitivity}, we can \textbf{lower the threshold}. Using 20\%:
\[
    \Pr(\text{default} = \text{yes} \mid X = x) > 0.2 \implies \text{predict Default}
\]

\begin{center}
\begin{tabular}{llccc}
\hline
 & & \multicolumn{3}{c}{True default status} \\
 & & No & Yes & Total \\
\hline
\multirow{2}{*}{Predicted default status} & No  & 9\,432 & 138 & 9\,570 \\
                                           & Yes & 235    & 195 & 430    \\
\hline
 & Total & 9\,667 & 333 & 10\,000 \\
\hline
\end{tabular}
\end{center}

With a 20\% threshold: sensitivity rises to $195/333 = 58.6\%$ (vs.\ 24.3\% before), but 235 non-defaulters are now incorrectly classified and the overall error rate increases slightly to 3.73\%.

\begin{tcolorbox}[colback=orange!5!white,colframe=orange!75!black,title=\textbf{Threshold Trade-off}]
Lowering the threshold \textbf{increases sensitivity} but \textbf{decreases specificity}. The optimal threshold depends on the \textbf{domain costs} associated with each error type and cannot be determined from the data alone — it requires external domain knowledge (e.g., the financial cost of a missed default vs.\ a false alarm).
\end{tcolorbox}

\subsection{The ROC Curve and AUC}

Rather than evaluating at a single threshold, the \textbf{ROC curve} (Receiver Operating Characteristics) simultaneously displays both error types \textbf{for all possible thresholds}:
\begin{itemize}
    \item $y$-axis: \textbf{True Positive Rate} (sensitivity).
    \item $x$-axis: \textbf{False Positive Rate} (1 $-$ specificity).
\end{itemize}

The \textbf{Area Under the ROC Curve (AUC)} summarises the overall classifier performance across all thresholds:
\begin{itemize}
    \item AUC $= 1$: perfect classifier.
    \item AUC $= 0.5$: no better than random chance (diagonal dotted line).
    \item The higher the AUC, the better the classifier.
\end{itemize}

For the LDA classifier on the Default data, AUC $= \mathbf{0.95}$ — excellent overall discrimination between defaulters and non-defaulters.

\begin{tcolorbox}[colback=green!5!white,colframe=green!75!black,title=\textbf{Key Takeaways for Classifier Evaluation}]
\begin{itemize}
    \item Never rely solely on the overall error rate: always inspect the \textbf{confusion matrix}.
    \item Report \textbf{sensitivity} and \textbf{specificity} separately — they reveal class-specific performance.
    \item Use the \textbf{ROC curve and AUC} to evaluate performance across all thresholds.
    \item Choose the decision threshold based on \textbf{domain knowledge} about the costs of each error type.
\end{itemize}
\end{tcolorbox}

\end{document}
